%% ──────────────────────────────────────────────
%%  فصل ۱۰ – مداخلات معاصر: پنج دهه‌ی اخیر
%%  فایل: chapters/ch10-contemporary.tex
%% ──────────────────────────────────────────────
\chapter{مداخلات معاصر: پنج دهه‌ی اخیر}
\label{ch:contemporary}

%% ── چکیده‌ی فصل ──
\begin{tcolorbox}[mybox, title={چکیده‌ی فصل}]
این فصل به بررسی مهم‌ترین مداخلات خارجی از دهه‌ی ۱۹۹۰ میلادی تا کنون می‌پردازد؛ دوره‌ای که فروپاشی نظم دوقطبی جنگ سرد، فضای مفهومی و عملیاتی تازه‌ای برای مداخله‌گری بین‌المللی گشود. از جنگ خلیج فارس (۱۹۹۱) و اشغال عراق (۲۰۰۳) گرفته تا فروپاشی سومالی، نسل‌کشی رواندا، بمباران کوزوو، بیست سال اشغال افغانستان، هرج‌ومرج لیبی پس از قذافی، فاجعه‌ی سوریه، جنگ یمن، و بحران اوکراین، هر یک الگویی متفاوت از تعامل میان \concept{نجات‌طلبی}، \concept{عمل‌گرایی}، و \concept{فرصت‌طلبی} را نمایان می‌سازند. در هر مورد، تحلیل \concept{شش‌وجهی} نشان می‌دهد که فاصله‌ی میان شعارهای انسان‌دوستانه و نتایج واقعی، معمولاً بسیار عمیق‌تر از آن است که ادبیات رسمی مداخله‌گران نشان می‌دهد.
\end{tcolorbox}

%% ================================================================
\section{درآمد: نظم نوین جهانی و بازتعریف مداخله}
\label{sec:ch10-intro}
%% ================================================================

\needspace{6\baselineskip}
با فروپاشی اتحاد جماهیر شوروی در سال ۱۹۹۱ و پایان جنگ سرد، ساختار دوقطبی‌ای که بیش از چهار دهه مانع یا محدودکننده‌ی بسیاری از مداخلات نظامی بود، از میان رفت. رئیس‌جمهور \histref{جرج بوش (پدر)}{۱۹۸۹–۱۹۹۳} در سخنرانی مشهور خود در کنگره‌ی آمریکا از \concept{نظم نوین جهانی} (\lat{New World Order}) سخن گفت؛ نظمی که بنا بود بر حاکمیت قانون بین‌المللی، حل‌وفصل مسالمت‌آمیز اختلافات، و اقدام جمعی علیه تجاوز استوار باشد.\footnote{\lr{Bush, George H.W., "Address Before a Joint Session of Congress on the State of the Union," January 29, 1991. \textit{Public Papers of the Presidents of the United States}, Government Printing Office, 1992.}}

اما همان‌گونه که در فصل‌های ۲ و ۳ دیدیم (بخش‌های \ref{ch:legal} و \ref{ch:models})، فاصله‌ی میان \concept{هنجار} و \concept{عمل} همواره فضای گسترده‌ای برای تفسیر، سوءاستفاده، و ابزارسازی فراهم می‌آورد. دهه‌ی ۱۹۹۰ و دو دهه‌ی نخست قرن بیست‌ویکم صحنه‌ی آزمایشگاهی بزرگی بود که در آن تقریباً هر نوع مداخله — از \concept{انسان‌دوستانه} تا \concept{رژیم‌چنج}، از \concept{حفظ صلح} تا \concept{اشغال}، و از \concept{تحریم} تا \concept{بمباران} — به اجرا درآمد، و نتایجی بسیار متنوع از رهایی نسبی تا فروپاشی کامل به‌بار آورد.

\begin{tcolorbox}[defbox, title={مفهوم کلیدی: مداخله‌ی معاصر}]
\concept{مداخله‌ی معاصر} (\lat{Contemporary Intervention}) به اَشکال گوناگون دخالت نظامی، سیاسی، اقتصادی، یا فرهنگی یک یا چند دولت (یا سازمان بین‌المللی) در امور داخلی دولت دیگر اطلاق می‌شود که پس از پایان جنگ سرد (۱۹۹۱) رخ داده و معمولاً با توجیهاتی چون حقوق بشر، مبارزه با تروریسم، یا مسئولیت حمایت (\lat{R2P}) همراه بوده است.
\end{tcolorbox}

\thinker{مایکل والزر} (\lat{Michael Walzer})، فیلسوف سیاسی آمریکایی، در ویرایش جدید کتاب \textit{جنگ‌های عادلانه و ناعادلانه} تأکید می‌کند که پایان جنگ سرد نه‌تنها مسئله‌ی مداخله را حل نکرد، بلکه آن را پیچیده‌تر ساخت؛ زیرا دیگر مانع بازدارنده‌ی اصلی — یعنی ترس از جنگ هسته‌ای — کاهش یافت و انگیزه‌های مداخله متنوع‌تر شدند.\footnote{\lr{Walzer, Michael. \textit{Just and Unjust Wars: A Moral Argument with Historical Illustrations}. 5th ed., New York: Basic Books, 2015, pp. 101–108.}}

در این فصل، ابتدا هر مورد مداخله‌ی مهم را به ترتیب زمانی بررسی می‌کنیم، سپس با استفاده از مدل \concept{شش‌وجهی نجات‌طلبی} (معرفی‌شده در فصل~\ref{ch:models}) حداقل سه مورد کلیدی (عراق ۲۰۰۳، لیبی ۲۰۱۱، و افغانستان ۲۰۰۱–۲۰۲۱) را تحلیل عمیق می‌کنیم، و در پایان جدول تطبیقی بزرگی از همه‌ی موارد ارائه می‌دهیم.


%% ================================================================
\section{عراق: از طوفان صحرا تا سقوط بغداد}
\label{sec:ch10-iraq}
%% ================================================================

\needspace{6\baselineskip}
\subsection{جنگ خلیج فارس ۱۹۹۱: مداخله‌ی «مشروع» و مرزهای آن}
\label{subsec:ch10-iraq-1991}

در ۲ اوت ۱۹۹۰، ارتش عراق به فرمان \histref{صدام حسین}{۱۹۳۷–۲۰۰۶} به کویت حمله کرد و طی چند ساعت آن را اشغال نمود. شورای امنیت سازمان ملل متحد با سرعتی بی‌سابقه قطعنامه‌ی ۶۶۰ را صادر کرد که خواستار خروج فوری عراق شد، و سپس با قطعنامه‌ی ۶۷۸ (۲۹ نوامبر ۱۹۹۰) به کشورهای عضو اجازه داد از «تمامی ابزارهای لازم» (\lat{all necessary means}) برای اجرای قطعنامه‌های پیشین استفاده کنند.\footnote{\lr{United Nations Security Council, Resolution 678 (1990), S/RES/678, 29 November 1990.}}

\begin{tcolorbox}[legalbox, title={نکته‌ی حقوقی: قطعنامه‌ی ۶۷۸}]
قطعنامه‌ی ۶۷۸ آخرین قطعنامه‌ی فصل هفتم شورای امنیت پیش از آغاز عملیات نظامی بود. عبارت «تمامی ابزارهای لازم» در عرف سازمان ملل معادل مجوز استفاده از زور تلقی می‌شود. این قطعنامه با ۱۲ رأی موافق، ۲ مخالف (کوبا و یمن)، و ۱ ممتنع (چین) تصویب شد.
\end{tcolorbox}

ائتلاف ۳۴ کشوری به رهبری ایالات متحده در ژانویه ۱۹۹۱ عملیات \concept{طوفان صحرا} (\lat{Desert Storm}) را آغاز کرد. پس از ۴۲ روز بمباران هوایی و ۱۰۰ ساعت عملیات زمینی، ارتش عراق از کویت بیرون رانده شد. اما \histref{جرج بوش (پدر)}{۱۹۸۹–۱۹۹۳} تصمیم گرفت عملیات را متوقف کند و به بغداد حمله نکند — تصمیمی که بعدها بسیار بحث‌برانگیز شد.\footnote{\lr{Gordon, Michael R. and Bernard E. Trainor. \textit{The Generals' War: The Inside Story of the Conflict in the Gulf}. Boston: Little, Brown, 1995, pp. 423–451.}}

مهم‌تر از خود جنگ، پیامدهای آن بود: شورش شیعیان جنوب و کردهای شمال عراق که با تشویق ضمنی آمریکا آغاز شد اما بدون حمایت واقعی رها شد. صدام حسین شورش‌ها را با خشونت سرکوب کرد. پاسخ جامعه‌ی بین‌المللی ایجاد \concept{مناطق پرواز ممنوع} (\lat{No-Fly Zones}) در شمال و جنوب عراق بود — اقدامی که مبنای حقوقی روشنی در قطعنامه‌های شورای امنیت نداشت اما بیش از یک دهه ادامه یافت.\footnote{\lr{Wheeler, Nicholas J. \textit{Saving Strangers: Humanitarian Intervention in International Society}. Oxford: Oxford University Press, 2000, pp. 139–171.}}

\begin{tcolorbox}[enemybox, title={فاجعه: سرکوب انتفاضه‌ی ۱۹۹۱}]
پس از آنکه رادیو \lat{Voice of Free Iraq} — وابسته به سیا — شیعیان و کردها را به قیام تشویق کرد، بوش تصمیم گرفت مداخله نکند. دلیل رسمی: «جلوگیری از بالکانیزه شدن عراق». برآوردها حاکی از کشته شدن ۳۰٬۰۰۰ تا ۶۰٬۰۰۰ نفر در سرکوب است. این رویداد نمونه‌ی بارز \concept{نجات‌طلبی ابزاری} بود: تشویق به قیام بدون تعهد به حمایت.\footnote{\lr{Makiya, Kanan. \textit{Cruelty and Silence: War, Tyranny, Uprising, and the Arab World}. New York: W.W. Norton, 1993, pp. 57–97.}}
\end{tcolorbox}

\subsection{تحریم‌ها و برنامه‌ی نفت در برابر غذا (۱۹۹۱–۲۰۰۳)}
\label{subsec:ch10-iraq-sanctions}

رژیم تحریم‌های سازمان ملل علیه عراق (قطعنامه‌ی ۶۶۱) یکی از گسترده‌ترین و طولانی‌ترین نظام‌های تحریمی تاریخ بود. برنامه‌ی \concept{نفت در برابر غذا} (\lat{Oil-for-Food Programme}) که از ۱۹۹۶ اجرا شد، قرار بود اثرات انسانی تحریم‌ها را کاهش دهد، اما خود به بزرگ‌ترین رسوایی فساد مالی سازمان ملل تبدیل شد.\footnote{\lr{Volcker, Paul A., Richard J. Goldstone, and Mark Pieth. \textit{The Management of the United Nations Oil-for-Food Programme}. Independent Inquiry Committee, 2005.}}

\thinker{جوی گوردون} (\lat{Joy Gordon}) در پژوهش جامع خود نشان داده است که تحریم‌ها طی دوازده سال منجر به مرگ صدها هزار نفر — به‌ویژه کودکان — شد، بدون آنکه تغییر معناداری در رفتار رژیم بعث ایجاد کند.\footnote{\lr{Gordon, Joy. \textit{Invisible War: The United States and the Iraq Sanctions}. Cambridge, MA: Harvard University Press, 2010, pp. 1–23.}} \thinker{دنیس هالیدی} (\lat{Denis Halliday})، هماهنگ‌کننده‌ی امور انسان‌دوستانه‌ی سازمان ملل در عراق، در سال ۱۹۹۸ از سمت خود استعفا داد و تحریم‌ها را \enemy{«نسل‌کشی»} خواند.\footnote{\lr{Halliday, Denis. "The Impact of the UN Sanctions on the People of Iraq," \textit{Journal of Palestine Studies}, Vol. 28, No. 2, 1999, pp. 29–37.}}

\subsection{اشغال عراق ۲۰۰۳: «نجات‌بخشی» به‌مثابه‌ی اشغال}
\label{subsec:ch10-iraq-2003}

حمله‌ی آمریکا و بریتانیا به عراق در ۲۰ مارس ۲۰۰۳ — بدون مجوز صریح شورای امنیت — نقطه‌ی عطفی در تاریخ مداخلات معاصر بود. توجیهات رسمی سه‌گانه بود: (۱) سلاح‌های کشتار جمعی، (۲) ارتباط با القاعده، و (۳) «آزادسازی» مردم عراق. هر سه توجیه بعدها بی‌اساس یا اغراق‌آمیز ثابت شدند.\footnote{\lr{Duelfer, Charles. \textit{Comprehensive Report of the Special Advisor to the DCI on Iraq's WMD} (Duelfer Report). Central Intelligence Agency, 2004.}}

\begin{tcolorbox}[quotebox, title={نقل‌قول تاریخی}]
«ما به‌عنوان آزادکنندگان آمده‌ایم، نه اشغالگران.»\\
— \histref{دونالد رامسفلد}{وزیر دفاع آمریکا}، مارس ۲۰۰۳\footnote{\lr{Ricks, Thomas E. \textit{Fiasco: The American Military Adventure in Iraq}. New York: Penguin Press, 2006, p. 79.}}
\end{tcolorbox}

واقعیت اما بسیار متفاوت بود. \thinker{توماس ریکس} (\lat{Thomas Ricks})، خبرنگار نظامی \textit{واشنگتن پست}، در کتاب \textit{فیاسکو} مستند کرده است که برنامه‌ریزی برای پس از جنگ تقریباً وجود نداشت. انحلال ارتش عراق به فرمان \histref{پل برمر}{مدیر اقتدار موقت ائتلاف} صدها هزار مرد مسلح بیکار و خشمگین ایجاد کرد که بسیاری‌شان به شورشیان پیوستند.\footnote{\lr{Bremer, L. Paul III. \textit{My Year in Iraq: The Struggle to Build a Future of Hope}. New York: Simon \& Schuster, 2006.}} سیاست \concept{بعث‌زدایی} (\lat{De-Ba'athification}) نیز ساختار اداری کشور را فلج کرد.\footnote{\lr{Allawi, Ali A. \textit{The Occupation of Iraq: Winning the War, Losing the Peace}. New Haven: Yale University Press, 2007, pp. 147–168.}}

\begin{tcolorbox}[failbox, title={عوامل شکست: دولت‌سازی در عراق}]
\begin{enumerate}
\item \textbf{فقدان برنامه‌ی پساجنگ:} بخش «آینده‌ی عراق» وزارت خارجه آمریکا نادیده گرفته شد
\item \textbf{انحلال ارتش:} فرمان شماره‌ی ۲ برمر — ۴۰۰٬۰۰۰ نظامی بیکار شدند
\item \textbf{بعث‌زدایی افراطی:} حذف طبقه‌ی مدیریتی کشور
\item \textbf{نادیده گرفتن ساختار قبیله‌ای و فرقه‌ای:} مدل دموکراسی لیبرال بر جامعه‌ی چندپاره تحمیل شد
\item \textbf{فساد سیستماتیک:} بازسازی به شرکت‌های آمریکایی (هالیبرتون) سپرده شد
\item \textbf{شکنجه و زندان ابوغریب:} مشروعیت «آزادکنندگان» را نابود کرد
\end{enumerate}
\end{tcolorbox}

هزینه‌ی انسانی اشغال عراق فاجعه‌بار بود. مطالعه‌ی مجله‌ی \textit{لنست} در سال ۲۰۰۶ تعداد مرگ‌های اضافی ناشی از جنگ را حدود ۶۵۵٬۰۰۰ نفر تخمین زد.\footnote{\lr{Burnham, Gilbert, et al. "Mortality after the 2003 invasion of Iraq: a cross-sectional cluster sample survey," \textit{The Lancet}, Vol. 368, No. 9545, 2006, pp. 1421–1428.}} پروژه‌ی \lat{Iraq Body Count} تا سال ۲۰۲۳ بیش از ۲۰۰٬۰۰۰ غیرنظامی کشته‌شده را مستند کرده است.\footnote{\lr{Iraq Body Count, \url{https://www.iraqbodycount.org/}, accessed 2024.}} هزینه‌ی مالی جنگ نیز طبق برآورد \thinker{جوزف استیگلیتز} (\lat{Joseph Stiglitz}) و \thinker{لیندا بیلمز} (\lat{Linda Bilmes}) بالغ بر ۳ تریلیون دلار شد.\footnote{\lr{Stiglitz, Joseph E. and Linda J. Bilmes. \textit{The Three Trillion Dollar War: The True Cost of the Iraq Conflict}. New York: W.W. Norton, 2008.}}

%% ── تحلیل شش‌وجهی عراق ──
\needspace{8\baselineskip}
\subsection{تحلیل شش‌وجهی: عراق ۲۰۰۳}
\label{subsec:ch10-iraq-hexagonal}

\begin{tcolorbox}[casebox, title={مطالعه‌ی موردی: عراق ۲۰۰۳ — تحلیل شش‌وجهی}]

\textbf{۱. بُعد ساختاری-سیاسی:}
عراق تحت حکومت حزب بعث و صدام حسین یک \concept{دولت اقتدارگرای سکولار} بود با ساختار قبیله‌ای-فرقه‌ای پیچیده (سنی، شیعه، کرد). رژیم اگرچه سرکوبگر بود، ثبات نسبی داشت. بحران مشروعیت اصلی از بیرون تحمیل شد — نه از درون.

\textbf{۲. بُعد حقوقی-هنجاری:}
هیچ قطعنامه‌ی صریحی برای استفاده از زور در ۲۰۰۳ صادر نشد. آمریکا و بریتانیا به قطعنامه‌ی ۱۴۴۱ (۲۰۰۲) استناد کردند، اما اکثر حقوقدانان بین‌المللی — از جمله \thinker{کوفی عنان}، دبیرکل سازمان ملل — جنگ را «غیرقانونی» دانستند.\footnote{\lr{Annan, Kofi. Interview with BBC, 16 September 2004: "I have indicated it was not in conformity with the UN charter."}}

\textbf{۳. بُعد اقتصادی:}
عراق دارای دومین ذخایر نفتی جهان بود. ارتباط میان نفت و تصمیم به حمله، اگرچه از سوی دولت آمریکا انکار شد، مستندات بعدی نشان داد که شرکت‌های نفتی در برنامه‌ریزی پساجنگ نقش داشتند.\footnote{\lr{Muttitt, Greg. \textit{Fuel on the Fire: Oil and Politics in Occupied Iraq}. London: Bodley Head, 2011.}}

\textbf{۴. بُعد نظامی-امنیتی:}
عدم توازن نظامی کامل: ارتش آمریکا ظرف سه هفته بغداد را فتح کرد. اما جنگ نامتقارن پس از اشغال هشت سال طول کشید و بیش از ۴٬۵۰۰ سرباز آمریکایی کشته شدند.

\textbf{۵. بُعد روان‌شناختی-فرهنگی:}
ترومای ۱۱ سپتامبر ۲۰۰۱ فضای روانی آمریکا را برای پذیرش جنگ پیشگیرانه آماده کرده بود. \concept{سوگیری تأیید} (\lat{Confirmation Bias}) در تفسیر اطلاعات استخباراتی نقش کلیدی داشت. در عراق نیز بخشی از مردم واقعاً از سقوط صدام خوشحال بودند، اما اکثریت نسبت به اشغال بیگانه مقاومت نشان دادند.

\textbf{۶. بُعد ژئوپولیتیک:}
\concept{نومحافظه‌کاران} (\lat{Neoconservatives}) آمریکایی — از جمله \histref{پل ولفوویتز}{معاون وزیر دفاع} و \histref{داگلاس فیث}{معاون سیاست دفاعی} — عراق را نقطه‌ی شروع بازسازی خاورمیانه‌ی بزرگ می‌دانستند. اسرائیل نیز از حذف صدام حمایت می‌کرد. ایران از سقوط دشمن دیرینه‌اش بهره برد اما با حضور نظامی آمریکا در همسایگی مواجه شد.\footnote{\lr{Mearsheimer, John J. and Stephen M. Walt. \textit{The Israel Lobby and U.S. Foreign Policy}. New York: Farrar, Straus and Giroux, 2007, pp. 229–262.}}
\end{tcolorbox}

\textbf{نتیجه بر طیف:} عراق ۲۰۰۳ در طیف نتایج (تعریف‌شده در فصل~\ref{ch:models}) بین \keyword{D} (رژیم دست‌نشانده و وابستگی ساختاری) و \keyword{C} (تغییر رژیم اما بی‌ثباتی) قرار می‌گیرد. دولت‌های پس از ۲۰۰۳ اگرچه منتخب بودند، اما عمیقاً فرقه‌ای، فاسد، و تحت نفوذ همزمان ایران و آمریکا عمل کردند.


%% ================================================================
\section{سومالی ۱۹۹۲–۱۹۹۵: شکست نخستین آزمایش}
\label{sec:ch10-somalia}
%% ================================================================

\needspace{6\baselineskip}
فروپاشی دولت سومالی پس از سقوط \histref{سیاد بره}{رئیس‌جمهور سومالی، ۱۹۶۹–۱۹۹۱} در ژانویه ۱۹۹۱، جنگ داخلی و قحطی گسترده‌ای را به‌دنبال آورد. تصاویر کودکان گرسنه‌ی سومالیایی در رسانه‌های غربی، فشار افکار عمومی برای مداخله را افزایش داد — پدیده‌ای که \thinker{پیرا رابینسون} (\lat{Piers Robinson}) آن را \concept{اثر سی‌ان‌ان} (\lat{CNN Effect}) نامید.\footnote{\lr{Robinson, Piers. \textit{The CNN Effect: The Myth of News, Foreign Policy and Intervention}. London: Routledge, 2002.}}

شورای امنیت با قطعنامه‌ی ۷۹۴ (دسامبر ۱۹۹۲) عملیات \concept{بازگرداندن امید} (\lat{Restore Hope}) را مجاز کرد — نخستین مورد مجوز استفاده از زور برای اهداف صرفاً انسان‌دوستانه بر اساس فصل هفتم.\footnote{\lr{United Nations Security Council, Resolution 794 (1992), S/RES/794, 3 December 1992.}} نیروهای آمریکایی و سپس یونیتف (\lat{UNITAF}) با حدود ۳۷٬۰۰۰ نیرو وارد شدند.

اما عملیات به‌سرعت از توزیع کمک‌های انسانی به تلاش برای خلع‌سلاح جنگ‌سالاران — به‌ویژه \histref{محمد فارح عیدید}{جنگ‌سالار}  — تغییر ماهیت داد. در ۳ اکتبر ۱۹۹۳، عملیاتی برای دستگیری معاونان عیدید در موگادیشو به فاجعه انجامید: دو هلی‌کوپتر بلک‌هاوک ساقط شدند، ۱۸ سرباز آمریکایی کشته و اجسادشان در خیابان‌ها کشیده شدند.\footnote{\lr{Bowden, Mark. \textit{Black Hawk Down: A Story of Modern War}. New York: Atlantic Monthly Press, 1999.}}

\begin{tcolorbox}[enemybox, title={نقطه‌ی عطف: موگادیشو ۱۹۹۳}]
تصاویر کشیده شدن جسد سربازان آمریکایی در خیابان‌های موگادیشو، \concept{سندرم سومالی} (\lat{Somalia Syndrome}) را در سیاست خارجی آمریکا ایجاد کرد: اکراه شدید از مداخله‌ی زمینی در مناطق بحرانی. این سندرم مستقیماً بر عدم مداخله در رواندا (۱۹۹۴) تأثیر گذاشت.
\end{tcolorbox}

\thinker{الکس دو وال} (\lat{Alex de Waal}) استدلال کرده است که مداخله در سومالی از ابتدا بر درک نادرست استوار بود: جامعه‌ی بین‌المللی سومالی را «دولت شکست‌خورده»ای تلقی کرد که به کمک بیرونی نیاز دارد، در حالی که ساختارهای سنتی حل اختلاف (قبیله‌ای و عرفی) وجود داشت که نادیده گرفته شد.\footnote{\lr{de Waal, Alex. "US War Crimes in Somalia," \textit{New Left Review}, No. 230, 1998, pp. 131–144.}}

سومالی پس از خروج نیروهای بین‌المللی در ۱۹۹۵، در طیف \keyword{E} (هرج‌ومرج و دولت شکست‌خورده) باقی ماند و تا دهه‌ها بدون دولت مرکزی مؤثر ادامه داد.


%% ================================================================
\section{رواندا ۱۹۹۴: قیمت عدم مداخله}
\label{sec:ch10-rwanda}
%% ================================================================

\needspace{6\baselineskip}
اگر سومالی نمونه‌ی مداخله‌ی شکست‌خورده بود، رواندا نمونه‌ی عدم مداخله‌ی فاجعه‌بار بود. طی حدود ۱۰۰ روز از آوریل تا ژوئیه ۱۹۹۴، بین ۸۰۰٬۰۰۰ تا ۱٬۰۰۰٬۰۰۰ نفر از توتسی‌ها و هوتوهای میانه‌رو به دست شبه‌نظامیان هوتو (\concept{اینتراهاموه}، \lat{Interahamwe}) قتل‌عام شدند.\footnote{\lr{Des Forges, Alison. \textit{Leave None to Tell the Story: Genocide in Rwanda}. New York: Human Rights Watch, 1999.}}

ژنرال \histref{رومئو دالر}{فرمانده‌ی یونامیر} (\lat{UNAMIR})، فرمانده‌ی نیروی حافظ صلح سازمان ملل در رواندا، بارها به نیویورک هشدار داد که نسل‌کشی در حال وقوع است و درخواست تقویت نیرو کرد، اما سازمان ملل — تحت فشار آمریکا — نه‌تنها نیرو نفرستاد، بلکه تعداد نیروهای موجود را کاهش داد.\footnote{\lr{Dallaire, Roméo. \textit{Shake Hands with the Devil: The Failure of Humanity in Rwanda}. Toronto: Random House Canada, 2003, pp. 221–265.}}

\begin{tcolorbox}[quotebox, title={نقل‌قول تاریخی}]
«من دست شیطان را فشردم. من آن را دیدم. بوی آن را استشمام کردم. و دنیا تماشا کرد.»\\
— ژنرال \histref{رومئو دالر}{فرمانده‌ی یونامیر، رواندا}\footnote{\lr{Dallaire, Roméo. \textit{Shake Hands with the Devil}, op. cit., p. xv.}}
\end{tcolorbox}

\thinker{سامانتا پاور} (\lat{Samantha Power})، در کتاب برنده‌ی جایزه‌ی پولیتزر خود، \textit{مشکلی از جهنم}، نشان داده است که دولت کلینتون عمداً از به‌کار بردن واژه‌ی «نسل‌کشی» خودداری کرد تا از تعهد حقوقی به مداخله (بر اساس کنوانسیون ۱۹۴۸) اجتناب کند.\footnote{\lr{Power, Samantha. \textit{A Problem from Hell: America and the Age of Genocide}. New York: Basic Books, 2002, pp. 329–389.}}

نقش فرانسه نیز بسیار بحث‌برانگیز بود. فرانسه از رژیم هوتوی \histref{ژوونال هابیاریمانا}{رئیس‌جمهور رواندا} حمایت می‌کرد و عملیات \concept{فیروزه} (\lat{Opération Turquoise}) در ژوئن ۱۹۹۴ — که ظاهراً انسان‌دوستانه بود — در عمل به فرار عاملان نسل‌کشی به زئیر (کنگو) کمک کرد.\footnote{\lr{Prunier, Gérard. \textit{The Rwanda Crisis: History of a Genocide}. New York: Columbia University Press, 1997, pp. 281–311.}}

\begin{tcolorbox}[wavebox, title={درس رواندا: پارادوکس مداخله}]
رواندا دو درس متناقض به جهان داد:
\begin{itemize}
\item \textbf{درس اول:} عدم مداخله می‌تواند فاجعه‌بار باشد — «تماشای نسل‌کشی» اخلاقاً غیرقابل‌دفاع است.
\item \textbf{درس دوم:} مداخله‌هایی که صورت گرفت (فرانسه) خود بخشی از مشکل بود — حمایت پیشین از رژیم عامل نسل‌کشی.
\end{itemize}
این پارادوکس مستقیماً به شکل‌گیری دکترین \concept{مسئولیت حمایت} (\lat{R2P}) در سال ۲۰۰۱ منجر شد (رجوع به فصل~\ref{ch:legal}).
\end{tcolorbox}

%% ================================================================
\section{کوزوو ۱۹۹۹: مداخله بدون مجوز}
\label{sec:ch10-kosovo}
%% ================================================================

\needspace{6\baselineskip}
بحران کوزوو ریشه در فروپاشی یوگسلاوی و سیاست‌های \histref{اسلوبودان میلوشویچ}{رئیس‌جمهور صربستان} داشت. پس از جنگ‌های بوسنی (۱۹۹۲–۱۹۹۵) و توافق دیتون، کوزوو — استانی خودمختار با اکثریت آلبانیایی‌تبار — به کانون بحران بعدی تبدیل شد. \concept{ارتش آزادی‌بخش کوزوو} (\lat{KLA/UÇK}) از ۱۹۹۶ مبارزه‌ی مسلحانه آغاز کرد و صربستان با سرکوب گسترده پاسخ داد.\footnote{\lr{Judah, Tim. \textit{Kosovo: War and Revenge}. 2nd ed., New Haven: Yale University Press, 2002, pp. 100–138.}}

ناتو در ۲۴ مارس ۱۹۹۹ بمباران صربستان را آغاز کرد — بدون مجوز شورای امنیت (به دلیل وتوی قطعی روسیه و چین). عملیات \concept{نیروی متحد} (\lat{Allied Force}) هفتادوهشت روز طول کشید و شامل بیش از ۳۸٬۰۰۰ پرواز جنگی بود.\footnote{\lr{Lambeth, Benjamin S. \textit{NATO's Air War for Kosovo: A Strategic and Operational Assessment}. Santa Monica: RAND Corporation, 2001.}}

\begin{tcolorbox}[legalbox, title={نکته‌ی حقوقی: مشروعیت بدون قانونیت؟}]
بمباران کوزوو بحث گسترده‌ای در حقوق بین‌الملل ایجاد کرد. کمیسیون مستقل بین‌المللی کوزوو نتیجه‌ی معروفی گرفت: مداخله‌ی ناتو \concept{«غیرقانونی اما مشروع»} (\lat{illegal but legitimate}) بود — یعنی مبنای حقوقی نداشت اما از نظر اخلاقی قابل دفاع بود.\footnote{\lr{Independent International Commission on Kosovo. \textit{The Kosovo Report: Conflict, International Response, Lessons Learned}. Oxford: Oxford University Press, 2000, p. 4.}}

این فرمول‌بندی، اگرچه از نظر سیاسی مفید بود، از نظر حقوقی خطرناک است: هر قدرتی می‌تواند مداخله‌ی خود را «غیرقانونی اما مشروع» بنامد.
\end{tcolorbox}

از نظر نتایج، کوزوو در نهایت به‌صورت یکجانبه در سال ۲۰۰۸ اعلام استقلال کرد — استقلالی که هنوز از سوی بسیاری از کشورها (از جمله روسیه، چین، اسپانیا، و صربستان) به رسمیت شناخته نمی‌شود. \thinker{نوام چامسکی} (\lat{Noam Chomsky}) استدلال کرده است که بمباران کوزوو نه‌تنها بحران انسانی را حل نکرد، بلکه در کوتاه‌مدت تشدید کرد: بیشتر \concept{پاک‌سازی قومی} (\lat{ethnic cleansing}) پس از آغاز بمباران رخ داد، نه قبل از آن.\footnote{\lr{Chomsky, Noam. \textit{The New Military Humanism: Lessons from Kosovo}. Monroe, ME: Common Courage Press, 1999, pp. 31–57.}}

روسیه از بحران کوزوو به‌شدت ناراحت شد و آن را نقض آشکار حاکمیت ملی و بی‌اعتنایی به شورای امنیت تلقی کرد. بعدها \histref{ولادیمیر پوتین}{رئیس‌جمهور روسیه} بارها به سابقه‌ی کوزوو برای توجیه الحاق کریمه در ۲۰۱۴ اشاره کرد.\footnote{\lr{Sakwa, Richard. \textit{Frontline Ukraine: Crisis in the Borderlands}. London: I.B. Tauris, 2015, pp. 74–91.}}

کوزوو در طیف نتایج بین \keyword{B} (رهایی با وابستگی محدود) و \keyword{C} (تغییر رژیم اما بی‌ثباتی) قرار می‌گیرد: مردم کوزوو از سرکوب صربستان رهایی یافتند، اما کشوری ناپایدار، فقیر، و وابسته به حضور نظامی ناتو (KFOR) باقی ماند.


%% ================================================================
\section{افغانستان ۲۰۰۱–۲۰۲۱: بیست سال اشغال و فروپاشی}
\label{sec:ch10-afghanistan}
%% ================================================================

\needspace{6\baselineskip}
\subsection{حمله و سقوط سریع طالبان}
\label{subsec:ch10-afghan-invasion}

حملات ۱۱ سپتامبر ۲۰۰۱ به مرکز تجارت جهانی و پنتاگون، دنیا را تغییر داد. دولت \histref{جرج دبلیو بوش}{رئیس‌جمهور آمریکا، ۲۰۰۱–۲۰۰۹} طالبان را به دلیل پناه دادن به \histref{اسامه بن‌لادن}{رهبر القاعده} تهدید کرد. شورای امنیت با قطعنامه‌ی ۱۳۶۸ (۱۲ سپتامبر ۲۰۰۱) حق دفاع مشروع را تأیید کرد، هرچند مجوز صریح حمله صادر نشد.\footnote{\lr{United Nations Security Council, Resolution 1368 (2001), S/RES/1368, 12 September 2001.}}

عملیات \concept{آزادی پایدار} (\lat{Enduring Freedom}) در ۷ اکتبر ۲۰۰۱ آغاز شد. نیروهای ویژه‌ی آمریکایی با همکاری \concept{ائتلاف شمال} (\lat{Northern Alliance}) ظرف دو ماه طالبان را از قدرت ساقطدند. اما — همان‌گونه که در عراق نیز دیدیم — پیروزی نظامی سریع به معنای موفقیت سیاسی نبود.

\subsection{دولت‌سازی شکست‌خورده}
\label{subsec:ch10-afghan-statebuilding}

کنفرانس بُن (دسامبر ۲۰۰۱) ساختار دولت موقت به ریاست \histref{حامد کرزی}{رئیس‌جمهور افغانستان، ۲۰۰۱–۲۰۱۴} را ایجاد کرد. اما پروژه‌ی دولت‌سازی با مشکلات بنیادین مواجه بود:

\begin{tcolorbox}[failbox, title={عوامل شکست: دولت‌سازی در افغانستان}]
\begin{enumerate}
\item \textbf{فقدان دولت مرکزی تاریخی:} افغانستان هرگز دولت-ملت متمرکزی نداشت
\item \textbf{ساختار قومی-قبیله‌ای:} پشتون، تاجیک، هزاره، ازبک — هر کدام منافع متفاوت
\item \textbf{جابه‌جایی توجه به عراق:} از ۲۰۰۳، منابع و توجه آمریکا به عراق منتقل شد
\item \textbf{فساد سیستماتیک:} دولت کرزی و اشرف‌غنی عمیقاً فاسد بودند
\item \textbf{پناهگاه امن پاکستان:} سازمان اطلاعات پاکستان (\lat{ISI}) از طالبان حمایت می‌کرد\footnote{\lr{Rashid, Ahmed. \textit{Pakistan on the Brink: The Future of America, Pakistan, and Afghanistan}. New York: Viking, 2012, pp. 15–44.}}
\item \textbf{اقتصاد مواد مخدر:} کشت خشخاش مهم‌ترین منبع درآمد بود
\item \textbf{تلفات غیرنظامی:} بمباران‌های هوایی و عملیات شبانه اعتماد مردم را نابود کرد
\end{enumerate}
\end{tcolorbox}

\thinker{سارا چیز} (\lat{Sarah Chayes}) که سال‌ها در قندهار زندگی کرد، در کتاب خود نشان داده است که فساد دولت کرزی خود عامل اصلی بازگشت طالبان بود: مردم نه از عشق به طالبان، بلکه از نفرت نسبت به فساد دولتی به آن‌ها روی آوردند.\footnote{\lr{Chayes, Sarah. \textit{Thieves of State: Why Corruption Threatens Global Security}. New York: W.W. Norton, 2015, pp. 57–83.}}

\subsection{خروج و فروپاشی ۲۰۲۱}
\label{subsec:ch10-afghan-withdrawal}

در ۱۵ اوت ۲۰۲۱، طالبان کابل را بدون جنگ تصرف کردند. ارتش ملی افغانستان — که ۲۰ سال و بیش از ۸۸ میلیارد دلار صرف آموزش و تجهیزش شده بود — طی چند هفته فروپاشید.\footnote{\lr{Special Inspector General for Afghanistan Reconstruction (SIGAR). \textit{What We Need to Learn: Lessons from Twenty Years of Afghanistan Reconstruction}. Washington, DC, 2021.}} صحنه‌های فرودگاه کابل — مردمی که از بال هواپیما آویزان بودند — نماد تلخ پایان بیست سال «مداخله‌ی آزادی‌بخش» شد.

گزارش بازرس ویژه‌ی آمریکا برای بازسازی افغانستان (\lat{SIGAR}) — معروف به «اوراق افغانستان» (\lat{Afghanistan Papers}) — نشان داد که مقامات آمریکایی سال‌ها می‌دانستند جنگ قابل پیروزی نیست اما به افکار عمومی دروغ می‌گفتند.\footnote{\lr{Whitlock, Craig. \textit{The Afghanistan Papers: A Secret History of the War}. New York: Simon \& Schuster, 2021.}}

%% ── تحلیل شش‌وجهی افغانستان ──
\needspace{8\baselineskip}
\subsection{تحلیل شش‌وجهی: افغانستان ۲۰۰۱–۲۰۲۱}
\label{subsec:ch10-afghan-hexagonal}

\begin{tcolorbox}[casebox, title={مطالعه‌ی موردی: افغانستان — تحلیل شش‌وجهی}]

\textbf{۱. بُعد ساختاری-سیاسی:}
افغانستان فاقد سنت دولت متمرکز بود. قدرت واقعی در دست جنگ‌سالاران، سران قبایل، و فرماندهان محلی بود. تلاش برای تحمیل مدل دولت-ملت لیبرال غربی بر این ساختار، محکوم به شکست بود.

\textbf{۲. بُعد حقوقی-هنجاری:}
مبنای حقوقی حمله نسبتاً قوی بود (دفاع مشروع پس از ۱۱ سپتامبر)، اما مبنای حقوقی اشغال بیست‌ساله و دولت‌سازی بسیار ضعیف‌تر بود.

\textbf{۳. بُعد اقتصادی:}
افغانستان منابع طبیعی قابل‌توجهی دارد (لیتیوم، مس، سنگ‌های قیمتی) اما فاقد زیرساخت استخراج بود. اقتصاد مبتنی بر کمک‌های خارجی و مواد مخدر بود. از مجموع ۲٫۳ تریلیون دلار هزینه‌ی آمریکا، بخش اعظم به پیمانکاران خود آمریکا بازگشت.\footnote{\lr{Costs of War Project, Watson Institute, Brown University, 2021, \url{https://watson.brown.edu/costsofwar/}.}}

\textbf{۴. بُعد نظامی-امنیتی:}
عدم توازن فناورانه‌ی مطلق (هواپیماهای بدون سرنشین، ماهواره‌ها) در برابر جنگ چریکی با کمین، بمب‌های کنار جاده‌ای، و حملات انتحاری. طالبان صبور بودند: «شما ساعت دارید، ما زمان داریم.»

\textbf{۵. بُعد روان‌شناختی-فرهنگی:}
\concept{پشتون‌والی} (قانون اخلاقی پشتون‌ها) مقاومت در برابر اشغالگر خارجی را ارزشی بنیادین می‌شمارد. تاریخ شکست‌های بریتانیا (سه جنگ انگلو-افغان) و شوروی (۱۹۷۹–۱۹۸۹) این هویت مقاومتی را تقویت کرده بود.

\textbf{۶. بُعد ژئوپولیتیک:}
افغانستان در چهارراه رقابت قدرت‌ها: آمریکا، پاکستان (حمایت از طالبان)، ایران (نفوذ در غرب)، هند (حمایت از دولت کابل)، چین (منابع معدنی)، و روسیه (رقابت تاریخی). هیچ‌یک از این بازیگران منفعتی واقعی در ثبات افغانستان نداشت.
\end{tcolorbox}

\textbf{نتیجه بر طیف:} افغانستان ۲۰۰۱–۲۰۲۱ نتیجه‌ای بین \keyword{E} (هرج‌ومرج و دولت شکست‌خورده) و \keyword{D} (رژیم دست‌نشانده) به‌بار آورد. با بازگشت طالبان، وضعیت عملاً به نقطه‌ی شروع بازگشت — با این تفاوت که ۲۰ سال، ده‌ها هزار کشته، و تریلیون‌ها دلار هزینه شده بود.


%% ================================================================
\section{لیبی ۲۰۱۱: مسئولیت حمایت و هرج‌ومرج}
\label{sec:ch10-libya}
%% ================================================================

\needspace{6\baselineskip}
\subsection{بهار عربی و قیام لیبی}
\label{subsec:ch10-libya-uprising}

در فوریه ۲۰۱۱، در پی خیزش‌های تونس و مصر، مردم \era{بنغازی} علیه حکومت چهل‌ودوساله‌ی \histref{معمر قذافی}{رهبر لیبی، ۱۹۶۹–۲۰۱۱} قیام کردند. قذافی با لحنی تهدیدآمیز اعلام کرد که شورشیان را «خانه‌به‌خانه، کوچه‌به‌کوچه» سرکوب خواهد کرد.\footnote{\lr{Chorin, Ethan. \textit{Exit the Colonel: The Hidden History of the Libyan Revolution}. New York: PublicAffairs, 2012, pp. 207–231.}}

شورای امنیت با قطعنامه‌ی ۱۹۷۳ (۱۷ مارس ۲۰۱۱) منطقه‌ی پرواز ممنوع اعلام کرد و به کشورهای عضو اجازه داد

```latex
«تمامی اقدامات لازم» برای حفاظت از غیرنظامیان انجام دهند — اما به‌صراحت «اشغال نظامی خارجی به هر شکل در هر بخش از خاک لیبی» را ممنوع کرد.\footnote{\lr{United Nations Security Council, Resolution 1973 (2011), S/RES/1973, 17 March 2011.}} روسیه و چین رأی ممتنع دادند اما وتو نکردند — تصمیمی که بعدها عمیقاً از آن پشیمان شدند.

\begin{tcolorbox}[legalbox, title={نکته‌ی حقوقی: قطعنامه‌ی ۱۹۷۳ و مسئولیت حمایت}]
قطعنامه‌ی ۱۹۷۳ نخستین بار بود که شورای امنیت به‌صراحت بر اساس دکترین \concept{مسئولیت حمایت} (\lat{Responsibility to Protect — R2P}) مجوز استفاده از زور علیه یک دولت عضو صادر کرد. از نظر حقوقی، این قطعنامه تنها حفاظت از غیرنظامیان را مجاز می‌کرد، نه تغییر رژیم. اما اجرای عملی بسیار فراتر رفت.
\end{tcolorbox}

\subsection{از حفاظت غیرنظامیان تا تغییر رژیم}
\label{subsec:ch10-libya-regimechange}

ناتو عملیات \concept{حامی یکپارچه} (\lat{Unified Protector}) را در ۳۱ مارس ۲۰۱۱ آغاز کرد. طی هفت ماه بمباران، بیش از ۹٬۷۰۰ پرواز حمله‌ای انجام شد.\footnote{\lr{NATO. "Operation Unified Protector: Final Mission Stats," 2 November 2011.}} آنچه به‌عنوان «حفاظت از غیرنظامیان» آغاز شده بود، به‌سرعت به حمایت هوایی مستقیم از شورشیان و تلاش آشکار برای سرنگونی قذافی تبدیل شد. نیروهای ویژه‌ی بریتانیا، فرانسه، و قطر به‌طور مخفیانه در خاک لیبی حضور داشتند و شورشیان را آموزش و هدایت می‌کردند.\footnote{\lr{Wehrey, Frederic. \textit{The Burning Shores: Inside the Battle for the New Libya}. New York: Farrar, Straus and Giroux, 2018, pp. 39–62.}}

در ۲۰ اکتبر ۲۰۱۱، قذافی در زادگاهش سِرت دستگیر و به‌طرز فجیعی کشته شد. \histref{هیلاری کلینتون}{وزیر خارجه‌ی آمریکا} در مصاحبه‌ای تلویزیونی با خنده گفت: «آمدیم، دیدیم، او مُرد» — بازنویسی جمله‌ی مشهور ژولیوس سزار.\footnote{\lr{Clinton, Hillary R. Interview with CBS News, 20 October 2011.}}

\begin{tcolorbox}[quotebox, title={نقل‌قول تاریخی}]
«هدف ما حمایت از غیرنظامیان بود، نه تغییر رژیم.»\\
— \histref{آندرس فوق راسموسن}{دبیرکل ناتو}، مارس ۲۰۱۱\\[6pt]
«من فکر نمی‌کنم این [لیبی] یک شکست بود... اشتباه، فقدان برنامه‌ریزی برای روز بعد بود.»\\
— \histref{باراک اوباما}{رئیس‌جمهور آمریکا}، آوریل ۲۰۱۶\footnote{\lr{Goldberg, Jeffrey. "The Obama Doctrine," \textit{The Atlantic}, April 2016.}}
\end{tcolorbox}

\subsection{لیبی پس از قذافی: دولت شکست‌خورده}
\label{subsec:ch10-libya-aftermath}

سقوط قذافی به جای دموکراسی، هرج‌ومرج آورد. لیبی به‌سرعت به میدان رقابت شبه‌نظامیان مسلح تبدیل شد. از ۲۰۱۴، دو دولت رقیب شکل گرفت: دولت \concept{توافق ملی} (\lat{GNA}) در طرابلس (مورد حمایت سازمان ملل و ترکیه) و مجلس نمایندگان در طبرق (مورد حمایت مصر، امارات، و روسیه) به فرماندهی ژنرال \histref{خلیفه حفتر}{فرمانده‌ی ارتش ملی لیبی}.\footnote{\lr{International Crisis Group. "Making the Best of France's Libya Summit," Middle East and North Africa Report No. 180, 2018.}}

\thinker{آلن کوپرمن} (\lat{Alan Kuperman}) در تحلیل مفصلی نشان داده است که روایت «قتل‌عام قریب‌الوقوع بنغازی» — که مبنای اصلی مداخله بود — اغراق‌آمیز بود. نیروهای قذافی در شهرهایی که قبلاً بازپس گرفته بودند (مثل اجدابیا) کشتار سیستماتیک غیرنظامیان انجام نداده بودند.\footnote{\lr{Kuperman, Alan J. "Obama's Libya Debacle: How a Well-Meaning Intervention Ended in Failure," \textit{Foreign Affairs}, Vol. 94, No. 2, 2015, pp. 66–77.}}

\begin{tcolorbox}[enemybox, title={فاجعه: پیامدهای مداخله در لیبی}]
\begin{itemize}
\item \textbf{جنگ داخلی مداوم:} لیبی از ۲۰۱۱ تاکنون در بی‌ثباتی به‌سر می‌برد
\item \textbf{بازار برده:} در سال ۲۰۱۷، CNN بازار برده‌فروشی مهاجران آفریقایی در لیبی را مستند کرد\footnote{\lr{Elbagir, Nima, et al. "People for Sale: Where Lives Are Auctioned for \$400," CNN, 14 November 2017.}}
\item \textbf{بحران مهاجرت:} لیبی به مسیر اصلی مهاجران غیرقانونی به اروپا تبدیل شد
\item \textbf{گسترش سلاح:} انبارهای عظیم تسلیحاتی قذافی به دست گروه‌های مسلح از مالی تا سوریه رسید
\item \textbf{ظهور داعش:} داعش توانست در سِرت (زادگاه قذافی) پایگاه ایجاد کند
\item \textbf{بی‌اعتمادی روسیه و چین:} هر دو کشور اعلام کردند ناتو از قطعنامه سوءاستفاده کرد — و این بی‌اعتمادی مستقیماً بر وتوهای بعدی درباره‌ی سوریه تأثیر گذاشت
\end{itemize}
\end{tcolorbox}

%% ── تحلیل شش‌وجهی لیبی ──
\needspace{8\baselineskip}
\subsection{تحلیل شش‌وجهی: لیبی ۲۰۱۱}
\label{subsec:ch10-libya-hexagonal}

\begin{tcolorbox}[casebox, title={مطالعه‌ی موردی: لیبی ۲۰۱۱ — تحلیل شش‌وجهی}]

\textbf{۱. بُعد ساختاری-سیاسی:}
لیبی تحت قذافی فاقد نهادهای سیاسی رسمی بود. \concept{جماهیریه} (\lat{Jamahiriya}) — «دولت توده‌ها» — در عمل حکومت فردی مطلق بود. قبایل نقش محوری در توزیع قدرت داشتند. سقوط قذافی، همه‌ی ساختارهای غیررسمی قدرت را نیز فروپاشاند.

\textbf{۲. بُعد حقوقی-هنجاری:}
مجوز شورای امنیت (قطعنامه ۱۹۷۳) وجود داشت، اما محدود به «حفاظت از غیرنظامیان» بود. تبدیل مأموریت به تغییر رژیم، نقض روح و احتمالاً حرف قطعنامه بود. \thinker{اریک پوزنر} (\lat{Eric Posner}) این را «چرخش مأموریت» (\lat{mission creep}) کلاسیک نامید.\footnote{\lr{Posner, Eric A. "Outside the Law: Two Views on the Libya Intervention," \textit{The New York Times}, 28 March 2011.}}

\textbf{۳. بُعد اقتصادی:}
لیبی بزرگ‌ترین ذخایر نفتی آفریقا را دارد (حدود ۴۸ میلیارد بشکه). نفت لیبی — سبک و کم‌گوگرد — برای پالایشگاه‌های اروپایی بسیار مطلوب بود. فرانسه‌ی \histref{نیکلا سارکوزی}{رئیس‌جمهور فرانسه} — نخستین کشور حامی مداخله — قراردادهای نفتی عمده‌ای با شورای انتقالی لیبی امضا کرد.\footnote{\lr{"Libyan Rebels Offered France 35\% of Oil in Exchange for Support," \textit{The Daily Telegraph}, 1 September 2011.}}

\textbf{۴. بُعد نظامی-امنیتی:}
عدم توازن مطلق: قذافی فاقد نیروی هوایی مؤثر و پدافند هوایی مدرن بود. بمباران ناتو ریسک تلفات نظامی اتحاد را نزدیک صفر نگه داشت — الگوی «جنگ بدون تلفات» (\lat{zero-casualty war}).

\textbf{۵. بُعد روان‌شناختی-فرهنگی:}
«بهار عربی» جوّ خوش‌بینانه‌ای در غرب ایجاد کرده بود: باور به اینکه ملت‌های عرب «بالاخره بیدار شده‌اند» و فقط نیاز به «کمکی از بیرون» دارند. این \concept{آرزواندیشی} (موضوع محوری این کتاب) واقعیت پیچیده‌ی جامعه‌ی لیبی را نادیده گرفت.

\textbf{۶. بُعد ژئوپولیتیک:}
فرانسه و بریتانیا انگیزه‌های متعدد داشتند: نفت، مبارزه با مهاجرت، رقابت با ایتالیا (متحد سنتی قذافی)، و اثبات قدرت نظامی پس از کاهش بودجه‌های دفاعی. آمریکا نقش «رهبری از پشت» (\lat{leading from behind}) را ایفا کرد — اصطلاحی که منتقدان آن را «عدم رهبری» خواندند.\footnote{\lr{Lizza, Ryan. "The Consequentialist: How the Arab Spring Remade Obama's Foreign Policy," \textit{The New Yorker}, 2 May 2011.}}
\end{tcolorbox}

\textbf{نتیجه بر طیف:} لیبی ۲۰۱۱ یکی از بارزترین نمونه‌های سقوط از \keyword{C} (تغییر رژیم اما بی‌ثباتی) به \keyword{E} (هرج‌ومرج و دولت شکست‌خورده) است. \histref{باراک اوباما}{رئیس‌جمهور آمریکا} خود بعدها لیبی را «بزرگ‌ترین اشتباه» دوران ریاست‌جمهوری‌اش نامید.\footnote{\lr{Goldberg, Jeffrey. "The Obama Doctrine," op. cit.}}


%% ================================================================
\section{سوریه ۲۰۱۱–کنون: فاجعه‌ی تمام‌عیار}
\label{sec:ch10-syria}
%% ================================================================

\needspace{6\baselineskip}
\subsection{از اعتراض تا جنگ داخلی}
\label{subsec:ch10-syria-civilwar}

در مارس ۲۰۱۱ — همزمان با موج «بهار عربی» — اعتراضات مسالمت‌آمیز در \era{درعا} علیه حکومت \histref{بشار اسد}{رئیس‌جمهور سوریه} آغاز شد. سرکوب خشونت‌آمیز رژیم، اعتراضات را به شورش مسلحانه و سپس به یکی از پیچیده‌ترین و ویرانگرترین جنگ‌های داخلی قرن بیست‌ویکم تبدیل کرد.\footnote{\lr{Lesch, David W. \textit{Syria: The Fall of the House of Assad}. New Haven: Yale University Press, 2012.}}

برخلاف لیبی، در مورد سوریه هیچ قطعنامه‌ی فصل هفتمی صادر نشد — روسیه و چین، که خود را فریفته‌ی قطعنامه‌ی ۱۹۷۳ (لیبی) می‌دانستند، بارها هرگونه اقدام الزام‌آور شورای امنیت را وتو کردند.\footnote{\lr{Zifcak, Spencer. "The Responsibility to Protect after Libya and Syria," \textit{Melbourne Journal of International Law}, Vol. 13, No. 1, 2012, pp. 59–93.}}

\subsection{بازیگران متعدد: سوریه به‌مثابه‌ی میدان جنگ نیابتی}
\label{subsec:ch10-syria-actors}

سوریه به میدان رقابت قدرت‌های منطقه‌ای و بین‌المللی تبدیل شد:

\begin{table}[htbp]
\centering
\caption{بازیگران اصلی جنگ سوریه و منافعشان}
\label{tab:ch10-syria-actors}
\begin{adjustbox}{max width=\linewidth}
\begin{tabularx}{\linewidth}{>{\raggedleft\arraybackslash}p{2.5cm} >{\raggedleft\arraybackslash}p{3cm} >{\raggedleft\arraybackslash}X}
\toprule
\rowcolor{PrimaryDeep}
\textcolor{white}{\textbf{بازیگر}} & \textcolor{white}{\textbf{طرف حمایت‌شده}} & \textcolor{white}{\textbf{انگیزه‌ی اصلی}} \\
\midrule
روسیه & دولت اسد & پایگاه دریایی طرطوس؛ حفظ نفوذ؛ مقابله با غرب \\
\rowcolor{NeutralLight}
ایران و حزب‌الله & دولت اسد & محور مقاومت؛ عمق استراتژیک \\
آمریکا & اپوزیسیون «میانه‌رو» و کردها & سرنگونی اسد؛ سپس مبارزه با داعش \\
\rowcolor{NeutralLight}
ترکیه & اپوزیسیون سنی و ارتش آزاد & سرنگونی اسد؛ مقابله با کردها \\
عربستان و قطر & گروه‌های اسلام‌گرا & مقابله با نفوذ ایران \\
\rowcolor{NeutralLight}
اسرائیل & — (حملات هوایی) & جلوگیری از انتقال سلاح به حزب‌الله \\
داعش & خودش & خلافت اسلامی \\
\bottomrule
\end{tabularx}
\end{adjustbox}
\end{table}

\thinker{کریستوفر فیلیپس} (\lat{Christopher Phillips}) در تحلیل جامع خود، جنگ سوریه را \concept{«نبرد برای سوریه»} نامیده و نشان داده است که بازیگران خارجی — بیش از بازیگران داخلی — مسیر جنگ را تعیین کردند.\footnote{\lr{Phillips, Christopher. \textit{The Battle for Syria: International Rivalry in the New Middle East}. New Haven: Yale University Press, 2016.}}

\subsection{سلاح شیمیایی و «خط قرمز» اوباما}
\label{subsec:ch10-syria-chemical}

در اوت ۲۰۱۳، حمله‌ی شیمیایی به غوطه‌ی شرقی (حومه‌ی دمشق) صدها نفر را کشت. اوباما پیش‌تر استفاده از سلاح شیمیایی را \concept{«خط قرمز»} اعلام کرده بود، اما در لحظه‌ی تصمیم‌گیری، از حمله‌ی نظامی عقب‌نشینی کرد و توافقی با روسیه برای انهدام سلاح‌های شیمیایی سوریه ترتیب داد.\footnote{\lr{Chollet, Derek. \textit{The Long Game: How Obama Defied Washington and Redefined America's Role in the World}. New York: PublicAffairs, 2016, pp. 149–178.}}

\begin{tcolorbox}[wavebox, title={تحلیل: پارادوکس خط قرمز}]
عقب‌نشینی اوباما از خط قرمز دو پیام متناقض ارسال کرد:
\begin{itemize}
\item \textbf{پیام اول:} آمریکا حاضر نیست برای هر بحرانی وارد جنگ شود — واقع‌گرایی.
\item \textbf{پیام دوم:} تهدیدهای آمریکا بی‌اعتبار است — ضعف.
\end{itemize}
\thinker{رابرت فورد} (\lat{Robert Ford})، آخرین سفیر آمریکا در دمشق، این تصمیم را «فاجعه‌بار» خواند: اپوزیسیون سوریه احساس خیانت کرد، روسیه جسورتر شد، و اسد فهمید که می‌تواند هر کاری بکند.\footnote{\lr{Ford, Robert. Testimony before the Senate Foreign Relations Committee, 2014.}}
\end{tcolorbox}

\subsection{مداخله‌ی نظامی روسیه (۲۰۱۵)}
\label{subsec:ch10-syria-russia}

در سپتامبر ۲۰۱۵، روسیه رسماً مداخله‌ی نظامی مستقیم در سوریه را آغاز کرد — به دعوت دولت اسد. این مداخله تعادل نظامی را به نفع اسد تغییر داد و عملاً بقای رژیم را تضمین کرد. بمباران‌های روسیه، که بخش قابل‌توجهی آن غیرنظامیان و بیمارستان‌ها را هدف قرار می‌داد، فاجعه‌ی انسانی را تشدید کرد.\footnote{\lr{Syrian Network for Human Rights. "Russia Has Killed More Syrian Civilians Than ISIS," Report, 2016.}}

حلب شرقی — که از ۲۰۱۲ تحت کنترل اپوزیسیون بود — پس از محاصره‌ی طولانی و بمباران بی‌رحمانه در دسامبر ۲۰۱۶ سقوط کرد. این رویداد نماد شکست جامعه‌ی بین‌المللی در حفاظت از غیرنظامیان شد.\footnote{\lr{Lund, Aron. "The Fall of Aleppo," \textit{The Century Foundation}, 14 December 2016.}}

\subsection{هزینه‌ی انسانی}
\label{subsec:ch10-syria-cost}

\begin{tcolorbox}[enemybox, title={فاجعه: آمار جنگ سوریه (تا ۲۰۲۳)}]
\begin{itemize}
\item \textbf{کشته‌شدگان:} بیش از ۵۰۰٬۰۰۰ نفر (برآورد دیدبان حقوق بشر سوریه)\footnote{\lr{Syrian Observatory for Human Rights, Cumulative Death Toll Report, 2023.}}
\item \textbf{آوارگان داخلی:} حدود ۶٫۹ میلیون نفر
\item \textbf{پناهجویان خارجی:} بیش از ۵٫۶ میلیون نفر (عمدتاً در ترکیه، لبنان، اردن)
\item \textbf{نیازمندان کمک انسانی:} ۱۴٫۶ میلیون نفر (از ۲۲ میلیون جمعیت پیش از جنگ)
\item \textbf{تخریب زیرساخت:} ۴۰۰ میلیارد دلار خسارت (برآورد بانک جهانی)\footnote{\lr{World Bank. \textit{The Toll of War: The Economic and Social Consequences of the Conflict in Syria}. 2017.}}
\end{itemize}
\end{tcolorbox}

سوریه از نظر طیف نتایج، مرکب از چند وضعیت همزمان است: مناطق تحت کنترل اسد (\keyword{D} — رژیم وابسته به روسیه و ایران)، شمال شرق کردنشین (\keyword{C} — خودمختاری ناپایدار وابسته به آمریکا)، شمال غربی اپوزیسیون (\keyword{E} — بی‌ثباتی تحت نفوذ ترکیه)، و مناطقی که داعش و بازماندگانش در آن فعال‌اند.


%% ================================================================
\section{یمن ۲۰۱۵–کنون: جنگ فراموش‌شده}
\label{sec:ch10-yemen}
%% ================================================================

\needspace{6\baselineskip}
بحران یمن ریشه‌ای پیچیده دارد: ترکیبی از تنش‌های فرقه‌ای (زیدی‌ها و سنی‌ها)، منطقه‌ای (شمال و جنوب)، قبیله‌ای، و ژئوپولیتیک (رقابت عربستان و ایران). پس از خیزش ۲۰۱۱ و سقوط \histref{علی عبدالله صالح}{رئیس‌جمهور یمن، ۱۹۷۸–۲۰۱۲}، دوره‌ی گذار تحت \histref{عبدربه منصور هادی}{رئیس‌جمهور انتقالی} شکست خورد و \concept{انصارالله} (حوثی‌ها) در سپتامبر ۲۰۱۴ صنعا را تصرف کردند.\footnote{\lr{Salisbury, Peter. "Yemen: National Chaos, Local Order," \textit{Chatham House Report}, December 2017.}}

در مارس ۲۰۱۵، ائتلاف به رهبری عربستان سعودی عملیات \concept{طوفان قاطعیت} (\lat{Decisive Storm}) را آغاز کرد — با هدف اعلام‌شده‌ی بازگرداندن دولت «مشروع» هادی. آمریکا، بریتانیا، و فرانسه با فروش تسلیحات، اطلاعات استخباراتی، و پشتیبانی لجستیکی از ائتلاف حمایت کردند.\footnote{\lr{Riedel, Bruce. "Who Are the Houthis, and Why Are We at War with Them?" \textit{Brookings Institution}, 18 December 2017.}}

\begin{tcolorbox}[enemybox, title={فاجعه: بحران انسانی یمن}]
سازمان ملل جنگ یمن را \enemy{«بدترین بحران انسانی جهان»} نامید:
\begin{itemize}
\item بیش از ۱۵۰٬۰۰۰ نفر کشته (تا ۲۰۲۱)
\item ۲۴٫۱ میلیون نفر (۸۰٪ جمعیت) نیازمند کمک انسانی
\item بزرگ‌ترین شیوع وبا در تاریخ معاصر: بیش از ۲٫۵ میلیون مورد مشکوک
\item محاصره‌ی بندر حدیده: قطع مسیر ۷۰٪ واردات غذا و دارو
\item حملات هوایی ائتلاف به مراسم عروسی، عزا، بیمارستان‌ها، و اتوبوس مدرسه\footnote{\lr{Mwatana for Human Rights. \textit{Day of Judgment: The Role of the US and Europe in Civilian Death, Destruction, and Trauma in Yemen}. 2019.}}
\end{itemize}
\end{tcolorbox}

از نظر حقوقی، مداخله‌ی عربستان مبنای بحث‌برانگیزی داشت: دعوت رئیس‌جمهور هادی — که خود مشروعیتش زیر سؤال بود — و \concept{حق دفاع جمعی} (ماده‌ی ۵۱ منشور). اما نبود قطعنامه‌ی صریح فصل هفتم و نقض گسترده‌ی حقوق بشردوستانه بین‌المللی، مشروعیت مداخله را عمیقاً تضعیف کرد.\footnote{\lr{Sharp, Jeremy M. "Yemen: Civil War and Regional Intervention," \textit{Congressional Research Service Report}, R43960, 2022.}}

\thinker{هلن لاکنر} (\lat{Helen Lackner})، یمن‌شناس برجسته، استدلال کرده است که مداخله‌ی عربستان نه‌تنها به اهداف اعلام‌شده‌اش نرسید، بلکه حوثی‌ها را قوی‌تر و محبوب‌تر کرد — نمونه‌ای دیگر از \concept{اثر بومرنگ} مداخلات خارجی.\footnote{\lr{Lackner, Helen. \textit{Yemen in Crisis: The Road to War}. London: Verso, 2019, pp. 238–264.}}


%% ================================================================
\section{اوکراین ۲۰۱۴–کنون: بازگشت جنگ بزرگ به اروپا}
\label{sec:ch10-ukraine}
%% ================================================================

\needspace{6\baselineskip}
\subsection{بحران کریمه و دونباس (۲۰۱۴)}
\label{subsec:ch10-ukraine-crimea}

بحران اوکراین نقطه‌ی عطفی در نظم پساجنگ سرد بود. پس از سقوط دولت \histref{ویکتور یانوکوویچ}{رئیس‌جمهور اوکراین} در فوریه ۲۰۱۴ (در پی اعتراضات \concept{میدان} / \lat{Euromaidan})، روسیه واکنش تهاجمی نشان داد: نیروهای بدون نشان نظامی — معروف به \concept{«آدم‌های سبز کوچولو»} (\lat{Little Green Men}) — شبه‌جزیره‌ی کریمه را اشغال کردند و در همه‌پرسی بحث‌برانگیزی آن را به روسیه ملحق کردند.\footnote{\lr{Sarotte, Mary Elise. \textit{Not One Inch: America, Russia, and the Making of Post-Cold War Stalemate}. New Haven: Yale University Press, 2021, pp. 398–421.}}

همزمان، در منطقه‌ی \concept{دونباس} (شرق اوکراین)، جدایی‌طلبان مورد حمایت روسیه «جمهوری‌های خلق» دونتسک و لوهانسک را اعلام کردند. جنگ کم‌شدتی آغاز شد که با توافقات \concept{مینسک ۱ و ۲} (۲۰۱۴ و ۲۰۱۵) ظاهراً مهار شد اما هرگز حل نشد.\footnote{\lr{Katchanovski, Ivan. "The Separatist War in Donbas: A Violent Break-up of Ukraine?" \textit{European Politics and Society}, Vol. 17, No. 4, 2016, pp. 473–489.}}

\begin{tcolorbox}[defbox, title={مفهوم کلیدی: جنگ ترکیبی}]
\concept{جنگ ترکیبی} (\lat{Hybrid Warfare}) اصطلاحی است که پس از بحران اوکراین رواج یافت و به ترکیب عملیات نظامی متعارف، جنگ سایبری، جنگ اطلاعاتی (اخبار جعلی)، عملیات پنهان، فشار اقتصادی، و استفاده از نیروهای نیابتی اشاره دارد. \thinker{والری گراسیموف} (\lat{Valery Gerasimov})، رئیس ستاد کل ارتش روسیه، در مقاله‌ی ۲۰۱۳ خود چارچوب نظری این رویکرد را ترسیم کرده بود.\footnote{\lr{Gerasimov, Valery. "The Value of Science Is in the Foresight," \textit{Military-Industrial Kurier}, 27 February 2013. English translation in \textit{Military Review}, January-February 2016.}}
\end{tcolorbox}

\subsection{تهاجم تمام‌عیار ۲۰۲۲}
\label{subsec:ch10-ukraine-2022}

در ۲۴ فوریه ۲۰۲۲، روسیه تهاجم تمام‌عیار به اوکراین را آغاز کرد — بزرگ‌ترین جنگ زمینی در اروپا از جنگ جهانی دوم. \histref{ولادیمیر پوتین}{رئیس‌جمهور روسیه} آن را \concept{«عملیات نظامی ویژه»} نامید و توجیهاتش شامل «غیرنازی‌سازی» و «غیرنظامی‌سازی» اوکراین، حفاظت از روس‌زبانان دونباس، و جلوگیری از عضویت اوکراین در ناتو بود.\footnote{\lr{Putin, Vladimir. "Address by the President of the Russian Federation," 24 February 2022. Official Kremlin transcript.}}

\begin{tcolorbox}[quotebox, title={نقل‌قول تاریخی}]
«من به حمل‌ونقل نیاز دارم، نه تاکسی.»\\
— \histref{ولودیمیر زلنسکی}{رئیس‌جمهور اوکراین}، در پاسخ به پیشنهاد آمریکا برای تخلیه‌اش از کی‌یف، فوریه ۲۰۲۲\footnote{\lr{Shuster, Simon. \textit{The Showman: Inside the Invasion That Shook the World and Made a Leader of Volodymyr Zelensky}. New York: William Morrow, 2024, p. 87.}}
\end{tcolorbox}

اوکراین با مقاومت سرسختانه، حمله‌ی اولیه‌ی روسیه به کی‌یف را دفع کرد. غرب (آمریکا، اتحادیه‌ی اروپا، بریتانیا) با ارسال تسلیحات پیشرفته، اطلاعات استخباراتی، تحریم‌های گسترده علیه روسیه، و حمایت مالی از اوکراین واکنش نشان داد — اما از مداخله‌ی نظامی مستقیم (به دلیل ترس از جنگ هسته‌ای) خودداری کرد.\footnote{\lr{Charap, Samuel and Timothy J. Colton. \textit{Everyone Loses: The Ukraine Crisis and the Ruinous Contest for Post-Soviet Eurasia}. London: Routledge, 2017; updated analysis in Charap, "An Unwinnable War," \textit{Foreign Affairs}, Vol. 102, No. 4, 2023.}}

\begin{tcolorbox}[wavebox, title={تحلیل: اوکراین و پارادوکس نجات‌طلبی}]
بحران اوکراین چند درس مهم دارد:
\begin{enumerate}
\item \textbf{مداخله‌ی غیرمستقیم:} غرب نشان داد که می‌توان بدون درگیری مستقیم، از طریق تسلیحات و تحریم‌ها مداخله کرد — اما این مداخله محدودیت‌های جدی دارد.
\item \textbf{بازگشت جنگ‌های بزرگ:} فرض اینکه جنگ‌های بزرگ بین‌دولتی پایان یافته، نادرست ثابت شد.
\item \textbf{نقش بازدارندگی هسته‌ای:} روسیه به‌عنوان قدرت هسته‌ای، سقف مداخله‌ی غرب را محدود کرد — بر خلاف عراق، لیبی، و افغانستان.
\item \textbf{هویت ملی و مقاومت:} مقاومت اوکراین نشان داد که عامل داخلی (اراده‌ی ملی) مهم‌تر از سلاح خارجی است.
\end{enumerate}
\end{tcolorbox}


%% ================================================================
\section{مداخلات منطقه‌ای آفریقایی: اوگاندا، کنگو، و الگوی «آفریقایی»}
\label{sec:ch10-africa}
%% ================================================================

\needspace{6\baselineskip}
آفریقا صحنه‌ی بسیاری از مداخلات منطقه‌ای بوده که کمتر مورد توجه رسانه‌های بین‌المللی قرار گرفته‌اند. دو مورد از مهم‌ترین آن‌ها عبارت‌اند از:

\subsection{جنگ‌های کنگو (۱۹۹۶–۲۰۰۳)}
\label{subsec:ch10-congo}

پیامد مستقیم نسل‌کشی رواندا، فرار عاملان هوتوی نسل‌کشی به \concept{زئیر} (\lat{Zaire}، کنگوی بعدی) بود. \histref{پل کاگامه}{رئیس‌جمهور رواندا} — به بهانه‌ی تعقیب عاملان نسل‌کشی — و \histref{یووری موسِوِنی}{رئیس‌جمهور اوگاندا} در ۱۹۹۶ به زئیر حمله کردند و رژیم \histref{موبوتو سسه سکو}{دیکتاتور زئیر} را سرنگون کردند. \histref{لوران-دزیره کبیلا}{رهبر شورشیان} به قدرت رسید.\footnote{\lr{Stearns, Jason K. \textit{Dancing in the Glory of Monsters: The Collapse of the Congo and the Great War of Africa}. New York: PublicAffairs, 2011.}}

اما جنگ دوم کنگو (۱۹۹۸–۲۰۰۳) — که «جنگ جهانی آفریقا» نامیده شده — نُه کشور آفریقایی را درگیر کرد و بیش از ۵٫۴ میلیون نفر تلفات (عمدتاً غیرمستقیم) داشت.\footnote{\lr{International Rescue Committee. "Mortality in the Democratic Republic of Congo: An Ongoing Crisis," 2007.}} این جنگ نمونه‌ای از \concept{مداخله‌ی منطقه‌ای} بود که در آن «نجات‌بخشی» (تعقیب عاملان نسل‌کشی) به غارت منابع (طلا، کلتان، الماس) تبدیل شد.

\subsection{اتحادیه‌ی آفریقا و مسئولیت مداخله}
\label{subsec:ch10-au}

یکی از تحولات مهم پساجنگ سرد، ایجاد \concept{اتحادیه‌ی آفریقا} (\lat{African Union — AU}) در ۲۰۰۲ و گنجاندن \concept{حق مداخله} در سند تأسیسی آن بود. ماده‌ی ۴(ح) اساسنامه‌ی اتحادیه‌ی آفریقا «حق اتحادیه برای مداخله در یک دولت عضو بنا بر تصمیم مجمع، در شرایط وخیم یعنی جنایات جنگی، نسل‌کشی، و جنایات علیه بشریت» را تصریح کرده است.\footnote{\lr{Constitutive Act of the African Union, Article 4(h), adopted in Lomé, Togo, 11 July 2000.}}

\begin{tcolorbox}[legalbox, title={نکته‌ی حقوقی: حق مداخله در اتحادیه‌ی آفریقا}]
اتحادیه‌ی آفریقا تنها سازمان منطقه‌ای جهان است که \concept{حق مداخله} در امور اعضا را به‌صراحت در سند تأسیسی‌اش گنجانده است — تحولی که مستقیماً از تجربه‌ی رواندا (عدم مداخله) و سومالی (شکست مداخله‌ی خارجی) برخاسته است. اما در عمل، این حق تقریباً هرگز بدون حمایت سازمان ملل و قدرت‌های بزرگ اِعمال نشده است.
\end{tcolorbox}

نمونه‌هایی از مداخلات آفریقایی شامل مأموریت اتحادیه‌ی آفریقا در سودان (\lat{AMIS})، مأموریت اتحادیه‌ی آفریقا در سومالی (\lat{AMISOM})، و مداخله‌ی اکووآس (\lat{ECOWAS}) در لیبریا و سیرالئون می‌شود. این مداخلات اگرچه نسبت به مداخلات غربی مشروعیت بیشتری داشته‌اند (به دلیل نزدیکی جغرافیایی و فرهنگی)، اما از نظر ظرفیت نظامی و مالی بسیار محدودتر بوده‌اند.\footnote{\lr{Williams, Paul D. \textit{War and Conflict in Africa}. 2nd ed., Cambridge: Polity Press, 2016, pp. 189–223.}}


%% ================================================================
\section{نمودار مقایسه‌ای: مداخلات معاصر بر اساس مدل شش‌وجهی}
\label{sec:ch10-tikz-comparison}
%% ================================================================

\needspace{12\baselineskip}

%% ── نمودار تیکز ۱: طیف نتایج ──
\begin{figure}[htbp]
\centering
\begin{tikzpicture}[
    >=Stealth,
    node distance=1.2cm,
    every node/.style={font=\small},
    scale=0.95, transform shape
]
% محور اصلی
\draw[->, thick, PrimaryDeep] (0,0) -- (14,0) node[below left]{\rl{بدترین نتیجه}};
\draw[->, thick, PrimaryDeep] (0,0) -- (0,8) node[above left, rotate=90, anchor=south]{\rl{بهترین نتیجه}};

% خطوط طیف
\foreach \y/\label/\grade in {
    7/{\rl{رهایی کامل}}/A,
    6/{\rl{رهایی با وابستگی}}/B,
    5/{\rl{تغییر رژیم + بی‌ثباتی}}/C,
    4/{\rl{رژیم دست‌نشانده}}/D,
    2/{\rl{هرج‌ومرج}}/E,
    1/{\rl{استعمار یا الحاق}}/F%
}{
    \draw[dashed, gray!40] (0.5,\y) -- (13.5,\y);
    \node[left] at (0,\y) {\textbf{\grade}};
    \node[right, gray] at (0.3,\y) {\label};
}

% موارد مداخله
\node[circle, fill=red!70, minimum size=8pt, inner sep=0pt, label={[font=\tiny]above:\rl{عراق ۲۰۰۳}}] at (2, 3.5) {};
\node[circle, fill=orange!70, minimum size=8pt, inner sep=0pt, label={[font=\tiny]above:\rl{لیبی ۲۰۱۱}}] at (4, 2.2) {};
\node[circle, fill=purple!70, minimum size=8pt, inner sep=0pt, label={[font=\tiny]below:\rl{افغانستان}}] at (6, 2) {};
\node[circle, fill=blue!60, minimum size=8pt, inner sep=0pt, label={[font=\tiny]above:\rl{کوزوو ۱۹۹۹}}] at (3, 5.5) {};
\node[circle, fill=black!70, minimum size=8pt, inner sep=0pt, label={[font=\tiny, white]above:\rl{سومالی}}] at (5, 1.5) {};
\node[circle, fill=brown!70, minimum size=8pt, inner sep=0pt, label={[font=\tiny]above:\rl{سوریه}}] at (8, 2.5) {};
\node[circle, fill=green!50!black, minimum size=8pt, inner sep=0pt, label={[font=\tiny]above:\rl{یمن}}] at (10, 2) {};
\node[circle, fill=cyan!70, minimum size=8pt, inner sep=0pt, label={[font=\tiny]above:\rl{اوکراین}}] at (12, 4.5) {};
\node[circle, fill=red!40, minimum size=8pt, inner sep=0pt, label={[font=\tiny]below:\rl{کریمه (روسیه)}}] at (11, 1) {};

% راهنما
\node[draw, rounded corners, fill=yellow!10, font=\tiny, text width=3.5cm, align=right] at (11.5, 7) {%
\rl{هر نقطه = یک مورد مداخله}\\%
\rl{محور عمودی = نتیجه بر طیف}\\%
\rl{محور افقی = ترتیب زمانی نسبی}%
};

\end{tikzpicture}
\caption{مقایسه‌ی نتایج مداخلات معاصر بر اساس طیف شش‌گانه}
\label{fig:ch10-spectrum}
\end{figure}


%% ── نمودار تیکز ۲: مدل علّی مداخلات شکست‌خورده ──
\needspace{14\baselineskip}

\begin{figure}[htbp]
\centering
\begin{tikzpicture}[
    >=Stealth,
    node distance=2cm,
    block/.style={rectangle, draw=PrimaryDeep, fill=PrimaryDeep!10, rounded corners,
        text width=3cm, minimum height=1.2cm, align=center, font=\small},
    cause/.style={rectangle, draw=red!70!black, fill=red!10, rounded corners,
        text width=2.8cm, minimum height=1cm, align=center, font=\footnotesize},
    result/.style={rectangle, draw=orange!70!black, fill=orange!10, rounded corners,
        text width=3cm, minimum height=1cm, align=center, font=\footnotesize},
    arrow/.style={->, thick, PrimaryDeep},
    scale=0.85, transform shape
]

% مرکز
\node[block] (intervention) {\rl{مداخله‌ی خارجی}};

% علل شکست (بالا)
\node[cause, above left=2.5cm and 1cm of intervention] (c1) {\rl{فقدان برنامه‌ی}\\%
\rl{پساجنگ}};
\node[cause, above=2.5cm of intervention] (c2) {\rl{نادیده گرفتن}\\%
\rl{ساختار اجتماعی}};
\node[cause, above right=2.5cm and 1cm of intervention] (c3) {\rl{آرزواندیشی و}\\%
\rl{سوگیری شناختی}};

% نتایج (پایین)
\node[result, below left=2.5cm and 1.5cm of intervention] (r1) {\rl{هرج‌ومرج}\\%
\rl{(سومالی، لیبی)}};
\node[result, below=2.5cm of intervention] (r2) {\rl{جنگ داخلی مزمن}\\%
\rl{(سوریه، یمن)}};
\node[result, below right=2.5cm and 1.5cm of intervention] (r3) {\rl{بازگشت به}\\%
\rl{نقطه‌ی صفر}\\%
\rl{(افغانستان)}};

% فلش‌ها
\draw[arrow] (c1) -- (intervention);
\draw[arrow] (c2) -- (intervention);
\draw[arrow] (c3) -- (intervention);
\draw[arrow, red!70!black] (intervention) -- (r1);
\draw[arrow, red!70!black] (intervention) -- (r2);
\draw[arrow, red!70!black] (intervention) -- (r3);

% علل فرعی
\node[cause, left=3cm of intervention] (c4) {\rl{منافع ژئوپولیتیک}\\%
\rl{بر انسان‌دوستی}};
\node[cause, right=3cm of intervention] (c5) {\rl{فقدان مشروعیت}\\%
\rl{حقوقی کافی}};
\draw[arrow] (c4) -- (intervention);
\draw[arrow] (c5) -- (intervention);

\end{tikzpicture}
\caption{مدل علّی شکست مداخلات معاصر}
\label{fig:ch10-causal-model}
\end{figure}


%% ================================================================
\section{جدول تطبیقی بزرگ: مداخلات معاصر}
\label{sec:ch10-comparative-table}
%% ================================================================

\needspace{6\baselineskip}
جدول زیر خلاصه‌ای تطبیقی از مهم‌ترین مداخلات معاصر بررسی‌شده در این فصل ارائه می‌دهد.

\begin{landscape}
\begin{table}[htbp]
\centering
\caption{جدول تطبیقی مداخلات معاصر (۱۹۹۱–۲۰۲۴)}
\label{tab:ch10-comparative}
\begin{adjustbox}{max width=\linewidth}
\begin{tabularx}{\linewidth}{
    >{\raggedleft\arraybackslash}p{1.8cm}
    >{\centering\arraybackslash}p{1.2cm}
    >{\raggedleft\arraybackslash}p{2.2cm}
    >{\raggedleft\arraybackslash}p{1.8cm}
    >{\centering\arraybackslash}p{1.2cm}
    >{\centering\arraybackslash}p{1.5cm}
    >{\raggedleft\arraybackslash}p{2.5cm}
    >{\raggedleft\arraybackslash}p{3cm}
    >{\centering\arraybackslash}p{1cm}
}
\toprule
\rowcolor{PrimaryDeep}
\textcolor{white}{\textbf{مورد}} &
\textcolor{white}{\textbf{سال}} &
\textcolor{white}{\textbf{مداخله‌گر}} &
\textcolor{white}{\textbf{نوع}} &
\textcolor{white}{\textbf{دعوت؟}} &
\textcolor{white}{\textbf{مجوز ش.ا.؟}} &
\textcolor{white}{\textbf{نتیجه‌ی فوری}} &
\textcolor{white}{\textbf{نتیجه‌ی بلندمدت}} &
\textcolor{white}{\textbf{طیف}} \\
\midrule

عراق (کویت) & ۱۹۹۱ & ائتلاف ۳۴ کشوری & نظامی & بلی & بلی & آزادی کویت & تحریم، مناطق پرواز ممنوع & B \\
\rowcolor{NeutralLight}
سومالی & ۱۹۹۲ & آمریکا/سازمان ملل & انسان‌دوستانه-نظامی & خیر & بلی & توزیع کمک & شکست، خروج، هرج‌ومرج & E \\

رواندا & ۱۹۹۴ & — (عدم مداخله) & عدم مداخله & — & — & نسل‌کشی & پیروزی RPF، ثبات اقتدارگرا & خاص \\
\rowcolor{NeutralLight}
کوزوو & ۱۹۹۹ & ناتو & نظامی (هوایی) & خیر & خیر & عقب‌نشینی صربستان & استقلال ۲۰۰۸، وابستگی & B/C \\

عراق & ۲۰۰۳ & آمریکا/بریتانیا & اشغال & خیر & خیر & سقوط صدام & جنگ فرقه‌ای، داعش & C/D \\
\rowcolor{NeutralLight}
افغانستان & ۲۰۰۱ & آمریکا/ناتو & اشغال & تا حدی & تا حدی & سقوط طالبان & بازگشت طالبان ۲۰۲۱ & E \\

لیبی & ۲۰۱۱ & ناتو & نظامی (هوایی) & تا حدی & بلی (R2P) & سقوط قذافی & هرج‌ومرج، دو دولت & E \\
\rowcolor{NeutralLight}
سوریه & ۲۰۱۱– & چندجانبه & نیابتی/مستقیم & متفاوت & خیر & جنگ داخلی & فاجعه‌ی انسانی & D/E \\

یمن & ۲۰۱۵– & عربستان/ائتلاف & نظامی & تا حدی & خیر & جنگ فرسایشی & بحران انسانی & E \\
\rowcolor{NeutralLight}
اوکراین (کریمه) & ۲۰۱۴ & روسیه & الحاق & خیر & خیر & الحاق کریمه & تنش مداوم & F \\

اوکراین & ۲۰۲۲– & روسیه/غرب (غیرمستقیم) & تهاجم/حمایت & خیر/بلی & خیر & مقاومت اوکراین & جنگ ادامه‌دار & ؟ \\
\rowcolor{NeutralLight}
کنگو & ۱۹۹۶– & رواندا/اوگاندا & منطقه‌ای & خیر & خیر & سقوط موبوتو & غارت منابع، ۵ میلیون کشته & E \\

\bottomrule
\end{tabularx}
\end{adjustbox}
\end{table}
\end{landscape}


%% ================================================================
\section{تحلیل تطبیقی

<!-- POSSIBLE OVERLAP DETECTED (similarity: 85%) - REVIEW BELOW -->

```latex
%% ================================================================
\section{تحلیل تطبیقی: الگوهای مشترک و متمایز}
\label{sec:ch10-comparative-analysis}
%% ================================================================

\needspace{6\baselineskip}
بررسی تطبیقی مداخلات معاصر — از جنگ خلیج فارس ۱۹۹۱ تا تهاجم روسیه به اوکراین در ۲۰۲۲ — چند الگوی مشترک و چند تمایز مهم را آشکار می‌سازد.

\subsection{الگوهای مشترک شکست}
\label{subsec:ch10-common-patterns}

\begin{tcolorbox}[wavebox, title={الگوی اول: شکاف میان توجیه و انگیزه}]
در تمامی موارد بررسی‌شده، \concept{شکاف معنادار}ی میان توجیهات رسمی (حقوق بشر، دموکراسی، مبارزه با تروریسم) و انگیزه‌های واقعی (نفت، ژئوپولیتیک، سیاست داخلی، رقابت قدرت‌ها) وجود داشته است. این شکاف هم مشروعیت مداخله را تضعیف می‌کند و هم برنامه‌ریزی پساجنگ را مختل می‌سازد — زیرا اهداف واقعی معمولاً با اهداف اعلام‌شده سازگار نیستند.

\textbf{مثال:} در عراق ۲۰۰۳، هدف اعلام‌شده «آزادسازی» بود، اما هدف واقعی بازسازی ژئوپولیتیک خاورمیانه بر اساس طرح نومحافظه‌کاران. در لیبی ۲۰۱۱، هدف اعلام‌شده «حفاظت از غیرنظامیان» بود، اما هدف واقعی تغییر رژیم و تأمین منافع نفتی. در اوکراین ۲۰۲۲ (از سوی روسیه)، هدف اعلام‌شده «غیرنازی‌سازی» بود، اما هدف واقعی بازسازی حوزه‌ی نفوذ شوروی سابق.
\end{tcolorbox}

\begin{tcolorbox}[wavebox, title={الگوی دوم: سندرم «پیروزی آسان، صلح غیرممکن»}]
در اکثر مداخلات نظامی معاصر، فاز نظامی با سرعت و سهولت نسبی انجام شده (سقوط صدام در سه هفته، سقوط طالبان در دو ماه، سقوط قذافی در هفت ماه)، اما فاز سیاسی — \concept{دولت‌سازی}، \concept{آشتی ملی}، و \concept{بازسازی} — یا اساساً برنامه‌ریزی نشده یا به‌شدت شکست خورده است.

\thinker{فرانسیس فوکویاما} (\lat{Francis Fukuyama})، که خود زمانی از مداخله در عراق حمایت کرده بود، بعدها اعتراف کرد که «مشکل اصلی این نبود که آیا می‌توان رژیمی را سرنگون کرد، بلکه این بود که آیا می‌توان چیز بهتری جایگزینش کرد.»\footnote{\lr{Fukuyama, Francis. \textit{State-Building: Governance and World Order in the 21st Century}. Ithaca: Cornell University Press, 2004, p. 99.}}
\end{tcolorbox}

\begin{tcolorbox}[wavebox, title={الگوی سوم: اثر بومرنگ}]
\concept{اثر بومرنگ} (\lat{Boomerang Effect}) به پدیده‌ای اشاره دارد که در آن مداخله‌ی خارجی نه‌تنها مشکل مورد نظر را حل نمی‌کند، بلکه مشکلات جدید و اغلب بدتری ایجاد می‌کند:
\begin{itemize}
\item \textbf{عراق ۲۰۰۳ → داعش:} اشغال عراق و بعث‌زدایی، زمینه‌ساز ظهور داعش در ۲۰۱۴ شد\footnote{\lr{Gerges, Fawaz A. \textit{ISIS: A History}. Princeton: Princeton University Press, 2016, pp. 89–127.}}
\item \textbf{لیبی ۲۰۱۱ → بحران ساحل:} سلاح‌های لیبی به دست گروه‌های مسلح در مالی، نیجر، و نیجریه رسید و بحران ساحل (\lat{Sahel Crisis}) را تشدید کرد\footnote{\lr{Lacher, Wolfram. "Libya's Fractious South and Regional Instability," \textit{Small Arms Survey}, HSBA Working Paper 37, 2014.}}
\item \textbf{افغانستان ۲۰۰۱ → بی‌ثباتی پاکستان:} جنگ افغانستان مناطق قبیله‌ای پاکستان را بی‌ثبات کرد و طالبان پاکستانی (\lat{TTP}) را تقویت نمود
\item \textbf{یمن ۲۰۱۵ → تقویت حوثی‌ها:} مداخله‌ی عربستان، حوثی‌ها را از یک گروه قبیله‌ای محلی به نیرویی منطقه‌ای با قابلیت حمله به زیرساخت‌های نفتی عربستان تبدیل کرد
\end{itemize}
\end{tcolorbox}

\subsection{الگوی چهارم: انتخابی بودن مداخله}
\label{subsec:ch10-selectivity}

یکی از مهم‌ترین انتقادات به نظام مداخله‌ی بین‌المللی، \concept{انتخابی بودن} (\lat{selectivity}) آن است. همان‌گونه که در فصل~\ref{ch:legal} بحث شد، اصل حاکمیت دولت‌ها در مواردی نقض و در موارد دیگر محترم شمرده می‌شود — نه بر اساس شدت نقض حقوق بشر، بلکه بر اساس منافع قدرت‌های بزرگ.

\begin{table}[htbp]
\centering
\caption{انتخابی بودن مداخلات: مقایسه‌ی واکنش‌ها}
\label{tab:ch10-selectivity}
\begin{adjustbox}{max width=\linewidth}
\begin{tabularx}{\linewidth}{>{\raggedleft\arraybackslash}p{2.5cm} >{\raggedleft\arraybackslash}p{3cm} >{\raggedleft\arraybackslash}p{3cm} >{\raggedleft\arraybackslash}X}
\toprule
\rowcolor{PrimaryDeep}
\textcolor{white}{\textbf{بحران}} &
\textcolor{white}{\textbf{شدت فاجعه}} &
\textcolor{white}{\textbf{واکنش بین‌المللی}} &
\textcolor{white}{\textbf{دلیل تفاوت واکنش}} \\
\midrule
رواندا ۱۹۹۴ & نسل‌کشی ۸۰۰٬۰۰۰ نفر & عدم مداخله & نبود منافع استراتژیک؛ سندرم سومالی \\
\rowcolor{NeutralLight}
کوزوو ۱۹۹۹ & پاک‌سازی قومی (ده‌ها هزار نفر) & بمباران ۷۸ روزه‌ی ناتو & اروپا؛ اعتبار ناتو؛ رسانه‌ها \\
لیبی ۲۰۱۱ & تهدید به سرکوب (نامشخص) & مداخله‌ی سریع ناتو & نفت؛ فرصت سیاسی؛ ضعف نظامی قذافی \\
\rowcolor{NeutralLight}
سوریه ۲۰۱۱– & ۵۰۰٬۰۰۰+ کشته & عدم مداخله‌ی مستقیم غرب & وتوی روسیه؛ پیچیدگی؛ سندرم عراق/لیبی \\
یمن ۲۰۱۵– & بدترین بحران انسانی & حمایت غرب از عربستان & فروش سلاح؛ رقابت با ایران \\
\rowcolor{NeutralLight}
میانمار ۲۰۱۷/۲۰۲۱ & نسل‌کشی روهینگیا؛ کودتا & تحریم‌های محدود & نبود منافع استراتژیک؛ نزدیکی به چین \\
فلسطین/غزه & بمباران مداوم غیرنظامیان & وتوی مکرر آمریکا & اتحاد استراتژیک آمریکا-اسرائیل \\
\rowcolor{NeutralLight}
چین (اویغورها) & بازداشت‌گاه‌های اجباری & اعتراض کلامی & قدرت اقتصادی و هسته‌ای چین \\
\bottomrule
\end{tabularx}
\end{adjustbox}
\end{table}

\thinker{ماروان بیشارا} (\lat{Marwan Bishara})، تحلیلگر الجزیره، به‌درستی اشاره کرده است که «حقوق بشر نه حق، بلکه \concept{امتیاز} است — امتیازی که تنها زمانی اِعمال می‌شود که با منافع قدرت‌ها همسو باشد.»\footnote{\lr{Bishara, Marwan. \textit{The Invisible Arab: The Promise and Peril of the Arab Revolutions}. New York: Nation Books, 2012, p. 203.}}

\subsection{الگوی پنجم: چرخش مأموریت}
\label{subsec:ch10-mission-creep}

\concept{چرخش مأموریت} (\lat{Mission Creep}) پدیده‌ای مکرر در مداخلات معاصر بوده است: مأموریتی که با هدف محدود آغاز می‌شود، به‌تدریج گسترش می‌یابد تا جایی که از هدف اولیه فاصله‌ی بنیادین پیدا می‌کند.

\begin{itemize}
\item \textbf{سومالی:} از توزیع کمک → به خلع‌سلاح جنگ‌سالاران → به جنگ شهری
\item \textbf{لیبی:} از حفاظت غیرنظامیان → به تغییر رژیم → به بمباران کاروان قذافی
\item \textbf{افغانستان:} از تعقیب بن‌لادن → به سرنگونی طالبان → به دولت‌سازی → به مبارزه با شورشیان → به مذاکره با طالبان
\end{itemize}

\thinker{رزا بروکس} (\lat{Rosa Brooks})، حقوقدان و مقام سابق پنتاگون، در کتاب خود نشان داده است که چرخش مأموریت ناشی از «ابهام ذاتی» مفهوم مداخله‌ی انسان‌دوستانه است: وقتی هدف «نجات مردم» باشد، تقریباً هر اقدامی قابل توجیه می‌شود.\footnote{\lr{Brooks, Rosa. \textit{How Everything Became War and the Military Became Everything}. New York: Simon \& Schuster, 2016, pp. 93–121.}}


%% ================================================================
\section{تفاوت‌های کلیدی: چرا برخی مداخلات «کمتر بد» بودند؟}
\label{sec:ch10-differences}
%% ================================================================

\needspace{6\baselineskip}
اگرچه اکثر مداخلات معاصر نتایج نامطلوبی داشته‌اند، برخی از آن‌ها نسبت به بقیه «کمتر بد» بوده‌اند. تحلیل تطبیقی نشان می‌دهد که عوامل زیر احتمال نتیجه‌ی نسبتاً بهتر را افزایش می‌دهند:

\begin{tcolorbox}[successbox, title={عوامل نسبی «موفقیت» در مداخلات}]
\begin{enumerate}
\item \textbf{اهداف محدود و مشخص:} مداخلاتی که هدف مشخص و محدود داشتند (مثلاً آزادسازی کویت ۱۹۹۱) نسبت به مداخلاتی با اهداف مبهم و گسترده (مثلاً «دولت‌سازی» در افغانستان) موفق‌تر بودند.

\item \textbf{مشروعیت حقوقی و بین‌المللی:} مداخلاتی با مجوز شورای امنیت و حمایت ائتلاف گسترده (کویت ۱۹۹۱) نسبت به مداخلات یکجانبه یا بدون مجوز (عراق ۲۰۰۳) نتایج بهتری داشتند.

\item \textbf{عامل داخلی قوی:} در کوزوو، \concept{ارتش آزادی‌بخش کوزوو} و ساختارهای مدنی آلبانیایی‌تبار زمینه‌ی داخلی برای دولت‌سازی فراهم کردند. در اوکراین، مقاومت ملی قوی عامل اصلی بقای کشور بود.

\item \textbf{نبود خلأ قدرت:} مداخلاتی که ساختارهای قدرت موجود را حفظ کردند (حتی با اصلاح) نسبت به مداخلاتی که همه‌چیز را نابود کردند (بعث‌زدایی در عراق، انحلال ساختارهای قذافی در لیبی) موفق‌تر بودند.

\item \textbf{تعهد بلندمدت واقعی:} در کوزوو، حضور KFOR و UNMIK بیش از دو دهه ادامه یافته و ثبات نسبی ایجاد کرده — بر خلاف سومالی که نیروها پس از شکست فوراً خارج شدند.
\end{enumerate}
\end{tcolorbox}

اما باید تأکید کرد که حتی «موفق‌ترین» مداخلات معاصر (کویت ۱۹۹۱ و تا حدی کوزوو ۱۹۹۹) نیز پیامدهای منفی قابل‌توجهی داشته‌اند. همان‌گونه که \thinker{اندرو بیسویچ} (\lat{Andrew Bacevich})، نظامی بازنشسته و استاد روابط بین‌الملل، نوشته است: «تاریخ مداخلات آمریکایی تاریخ امیدهای بزرگ و نتایج ناامیدکننده است.»\footnote{\lr{Bacevich, Andrew J. \textit{America's War for the Greater Middle East: A Military History}. New York: Random House, 2016, p. 358.}}


%% ================================================================
\section{نقش فناوری و رسانه در مداخلات معاصر}
\label{sec:ch10-tech-media}
%% ================================================================

\needspace{6\baselineskip}
یکی از ویژگی‌های متمایز مداخلات معاصر نسبت به دوره‌های پیشین (بررسی‌شده در فصل‌های ۸ و ۹)، نقش محوری فناوری و رسانه است.

\subsection{اثر سی‌ان‌ان و رسانه‌های ۲۴ ساعته}
\label{subsec:ch10-cnn-effect}

مفهوم \concept{اثر سی‌ان‌ان} (\lat{CNN Effect}) اشاره به تأثیر پوشش رسانه‌ای زنده بر تصمیم‌گیری سیاست خارجی دارد. تصاویر کودکان گرسنه‌ی سومالیایی (۱۹۹۲) مداخله را تسریع کرد؛ تصاویر اجساد سربازان آمریکایی در موگادیشو (۱۹۹۳) خروج را تسریع کرد. تصاویر پناهجویان کرد عراقی (۱۹۹۱) مناطق پرواز ممنوع را ممکن ساخت.\footnote{\lr{Livingston, Steven. "Clarifying the CNN Effect: An Examination of Media Effects According to Type of Military Intervention," \textit{Joan Shorenstein Center Research Paper}, R-18, Harvard University, 1997.}}

اما تحقیقات بعدی نشان داده‌اند که رابطه‌ی رسانه و سیاست خارجی پیچیده‌تر از آن است که «اثر سی‌ان‌ان» نشان می‌دهد. رسانه‌ها بیشتر در زمان \concept{ابهام سیاستی} (\lat{policy uncertainty}) تأثیرگذارند — یعنی زمانی که دولت هنوز تصمیم خود را نگرفته است.\footnote{\lr{Robinson, Piers. "The CNN Effect Revisited," \textit{Critical Studies in Media Communication}, Vol. 22, No. 4, 2005, pp. 344–349.}}

\subsection{شبکه‌های اجتماعی و «بهار عربی»}
\label{subsec:ch10-social-media}

ظهور شبکه‌های اجتماعی (فیسبوک، توییتر، یوتیوب) در دهه‌ی ۲۰۱۰ دینامیک مداخلات را تغییر داد. در «بهار عربی» ۲۰۱۱، شبکه‌های اجتماعی هم ابزار بسیج معترضان بودند و هم ابزار تبلیغاتی برای توجیه مداخله. ویدئوهای سرکوب خشونت‌آمیز در سوریه و لیبی از طریق یوتیوب جهانی شدند و فشار افکار عمومی برای مداخله را افزایش دادند.\footnote{\lr{Howard, Philip N. and Muzammil M. Hussain. \textit{Democracy's Fourth Wave? Digital Media and the Arab Spring}. Oxford: Oxford University Press, 2013.}}

\begin{tcolorbox}[defbox, title={مفهوم کلیدی: جنگ روایت‌ها}]
\concept{جنگ روایت‌ها} (\lat{Narrative Warfare}) به رقابت بازیگران مختلف برای تعریف «واقعیت» یک بحران اشاره دارد. در هر مداخله‌ی معاصر، چندین روایت رقیب وجود دارد:
\begin{itemize}
\item روایت مداخله‌گر: «ما برای نجات مردم آمده‌ایم»
\item روایت رژیم هدف: «این تجاوز به حاکمیت ماست»
\item روایت مخالفان داخلی: «ما از بیگانگان کمک خواستیم»
\item روایت قدرت‌های رقیب: «این نقض حقوق بین‌الملل است»
\end{itemize}
هر روایت بخشی از واقعیت را منعکس می‌کند، اما هیچ‌کدام تمام واقعیت نیست. تحلیل‌گر بی‌طرف باید بتواند همه‌ی روایت‌ها را بشنود و ارزیابی کند.
\end{tcolorbox}

\subsection{جنگ سایبری و هوش مصنوعی}
\label{subsec:ch10-cyber}

نسل جدید مداخلات شامل ابعاد سایبری نیز هست. حمله‌ی \concept{استاکس‌نت} (\lat{Stuxnet}) — ویروس رایانه‌ای آمریکایی-اسرائیلی — به تأسیسات هسته‌ای ایران در ۲۰۱۰، نخستین نمونه‌ی مستند «سلاح سایبری» بود که علیه زیرساخت فیزیکی یک کشور به‌کار رفت.\footnote{\lr{Zetter, Kim. \textit{Countdown to Zero Day: Stuxnet and the Launch of the World's First Digital Weapon}. New York: Crown, 2014.}} مداخله‌ی سایبری روسیه در انتخابات ۲۰۱۶ آمریکا و ۲۰۱۷ فرانسه، و استفاده‌ی گسترده از اطلاعات جعلی (\lat{disinformation}) در بحران اوکراین، نشان‌دهنده‌ی ابعاد جدید مداخله‌ی خارجی در عصر دیجیتال است.\footnote{\lr{Rid, Thomas. \textit{Active Measures: The Secret History of Disinformation and Political Warfare}. New York: Farrar, Straus and Giroux, 2020.}}

استفاده از \concept{پهپادها} (\lat{drones}) نیز ماهیت مداخله را تغییر داده است. برنامه‌ی کشتار هدفمند آمریکا (\lat{targeted killing}) با پهپاد در پاکستان، یمن، و سومالی شکل جدیدی از مداخله‌ی نظامی ایجاد کرد: مداخله‌ای بدون حضور نیروی زمینی، بدون ریسک تلفات خودی، اما با تلفات غیرنظامی قابل‌توجه.\footnote{\lr{Scahill, Jeremy. \textit{The Assassination Complex: Inside the Government's Secret Drone Warfare Program}. New York: Simon \& Schuster, 2016.}}


%% ================================================================
\section{نقد مسئولیت حمایت در پرتو تجربه‌ی معاصر}
\label{sec:ch10-r2p-critique}
%% ================================================================

\needspace{6\baselineskip}
دکترین \concept{مسئولیت حمایت} (\lat{R2P}) — که در فصل~\ref{ch:legal} به تفصیل بررسی شد — پس از تجربه‌های دهه‌ی ۲۰۱۰ با انتقادات جدی مواجه شده است.

\subsection{لیبی: اوج و سقوط R2P}
\label{subsec:ch10-r2p-libya}

لیبی ۲۰۱۱ هم اوج و هم سقوط R2P بود. از یک سو، قطعنامه‌ی ۱۹۷۳ نشان داد که شورای امنیت \textit{می‌تواند} بر اساس R2P عمل کند. از سوی دیگر، سوءاستفاده‌ی ناتو از قطعنامه (تبدیل حفاظت غیرنظامیان به تغییر رژیم) اعتماد روسیه و چین را نابود کرد و عملاً R2P را برای سال‌ها فلج ساخت.\footnote{\lr{Garwood-Gowers, Andrew. "The Responsibility to Protect and the Arab Spring: Libya as the Exception, Syria as the Norm?" \textit{UNSW Law Journal}, Vol. 36, No. 2, 2013, pp. 594–618.}}

\thinker{آلکس بلامی} (\lat{Alex Bellamy})، یکی از مهم‌ترین نظریه‌پردازان R2P، خود اذعان کرده است که «لیبی بیشتر از آنچه R2P را تقویت کرد، آن را تضعیف نمود.»\footnote{\lr{Bellamy, Alex J. \textit{The Responsibility to Protect: A Defense}. Oxford: Oxford University Press, 2015, p. 187.}}

\subsection{نقدهای ساختاری R2P}
\label{subsec:ch10-r2p-structural}

\begin{tcolorbox}[failbox, title={نقدهای بنیادین بر مسئولیت حمایت}]
\begin{enumerate}
\item \textbf{انتخابی بودن:} R2P در لیبی اعمال شد اما در سوریه، یمن، میانمار، و فلسطین نه — زیرا اجرایش به خواست اعضای دائم شورای امنیت وابسته است.

\item \textbf{سوءاستفاده‌پذیری:} هر قدرتی می‌تواند ادعا کند که برای «حفاظت از غیرنظامیان» مداخله می‌کند — حتی روسیه در اوکراین.\footnote{\lr{Hehir, Aidan. \textit{Hollow Norms and the Responsibility to Protect}. London: Palgrave Macmillan, 2019.}}

\item \textbf{نقض حاکمیت:} بسیاری از کشورهای جنوب جهانی (\lat{Global South}) — از جمله هند، برزیل، و آفریقای جنوبی — R2P را ابزار نوین سلطه‌ی غرب می‌دانند و مفهوم جایگزین \concept{«مسئولیت ضمن حمایت»} (\lat{Responsibility While Protecting — RwP}) را پیشنهاد کرده‌اند.\footnote{\lr{Benner, Thorsten. "Brazil as a Norm Entrepreneur: The 'Responsibility While Protecting' Initiative," \textit{Global Public Policy Institute Working Paper}, 2013.}}

\item \textbf{نبود مکانیزم نظارت:} پس از مجوز شورای امنیت، هیچ مکانیزم مؤثری برای نظارت بر عدم تجاوز از حدود مأموریت وجود ندارد.

\item \textbf{تضاد با واقعیت ژئوپولیتیک:} اجرای R2P علیه اعضای دائم شورای امنیت (مثلاً روسیه در چچن یا آمریکا در عراق) عملاً غیرممکن است.
\end{enumerate}
\end{tcolorbox}


%% ================================================================
\section{نمودار: تکامل دکترین مداخله از ۱۹۹۱ تا ۲۰۲۴}
\label{sec:ch10-tikz-evolution}
%% ================================================================

\needspace{14\baselineskip}

\begin{figure}[htbp]
\centering
\begin{tikzpicture}[
    >=Stealth,
    timeline/.style={very thick, PrimaryDeep},
    event/.style={circle, fill=PrimaryDeep, minimum size=6pt, inner sep=0pt},
    label above/.style={above=8pt, font=\footnotesize, text width=2.5cm, align=center},
    label below/.style={below=8pt, font=\footnotesize, text width=2.5cm, align=center},
    doctrine/.style={rectangle, draw=AccentGold, fill=AccentGold!15, rounded corners, font=\footnotesize, text width=2.8cm, align=center, minimum height=0.9cm},
    scale=0.75, transform shape
]

% خط زمان
\draw[timeline] (0,0) -- (20,0);

% رویدادها — بالا
\foreach \x/\yr/\txt in {
    1/1991/{\rl{جنگ خلیج فارس}},
    3/1992/{\rl{سومالی}},
    4.5/1994/{\rl{رواندا}},
    6.5/1999/{\rl{کوزوو}},
    8.5/2001/{\rl{افغانستان}},
    10.5/2003/{\rl{عراق}},
    13/2011/{\rl{لیبی}},
    15/2015/{\rl{یمن}},
    17.5/2022/{\rl{اوکراین}}%
}{
    \node[event] at (\x, 0) {};
    \node[label above] at (\x, 0) {\txt\\{\tiny \yr}};
}

% دکترین‌ها — پایین
\node[doctrine] at (2, -2) {\rl{نظم نوین جهانی}\\{\tiny 1991}};
\node[doctrine] at (5.5, -2) {\rl{سندرم سومالی}\\{\tiny 1993}};
\node[doctrine] at (8, -3.5) {\rl{مسئولیت حمایت}\\{\tiny \lr{R2P} — 2001}};
\node[doctrine] at (11, -2) {\rl{دکترین بوش}\\{\tiny \rl{جنگ پیشگیرانه}}};
\node[doctrine] at (14, -3.5) {\rl{رهبری از پشت}\\{\tiny \rl{اوباما}}};
\node[doctrine] at (17, -2) {\rl{بازگشت جنگ بزرگ}\\{\tiny \rl{رقابت قدرت‌ها}}};

% فلش‌ها از دکترین به خط زمان
\draw[->, dashed, gray] (2, -1.4) -- (1, -0.3);
\draw[->, dashed, gray] (5.5, -1.4) -- (4.5, -0.3);
\draw[->, dashed, gray] (8, -2.9) -- (6.5, -0.3);
\draw[->, dashed, gray] (11, -1.4) -- (10.5, -0.3);
\draw[->, dashed, gray] (14, -2.9) -- (13, -0.3);
\draw[->, dashed, gray] (17, -1.4) -- (17.5, -0.3);

% عنوان
\node[font=\normalsize\bfseries, PrimaryDeep] at (10, 2) {\rl{تکامل دکترین مداخله: ۱۹۹۱–۲۰۲۴}};

\end{tikzpicture}
\caption{خط زمانی تکامل دکترین‌های مداخله از ۱۹۹۱ تا ۲۰۲۴}
\label{fig:ch10-timeline}
\end{figure}


%% ================================================================
\section{مداخله‌ی معاصر و مسئله‌ی ایران}
\label{sec:ch10-iran-link}
%% ================================================================

\needspace{6\baselineskip}
اگرچه تحلیل مفصل وضعیت ایران پس از ۱۳۵۷ به پیوست~پ (\ref{app:iran-post1979}) واگذار شده، در اینجا لازم است تا به نقاط تماس میان مداخلات معاصر و ایران اشاره شود.

ایران در پنج دهه‌ی اخیر هم \textit{موضوع} مداخله و هم \textit{عامل} مداخله بوده است:

\begin{tcolorbox}[casebox, title={ایران: هم‌ موضوع، هم عامل مداخله}]
\textbf{ایران به‌عنوان موضوع مداخله:}
\begin{itemize}
\item تحریم‌های اقتصادی گسترده (از ۱۹۷۹ تاکنون)
\item حمله‌ی سایبری استاکس‌نت (۲۰۱۰)
\item ترور دانشمندان هسته‌ای (۲۰۱۰–۲۰۲۰)
\item فعالیت رسانه‌های فارسی‌زبان خارجی (BBC فارسی، صدای آمریکا، ایران اینترنشنال)
\item حمایت از جنبش‌های اپوزیسیون و گروه‌های تجزیه‌طلب
\item ترور ژنرال \histref{قاسم سلیمانی}{فرمانده‌ی نیروی قدس} توسط آمریکا (ژانویه ۲۰۲۰)\footnote{\lr{Filkins, Dexter. "The Shadow Commander," \textit{The New Yorker}, 30 September 2013; updated analysis post-assassination, January 2020.}}
\end{itemize}

\textbf{ایران به‌عنوان عامل مداخله:}
\begin{itemize}
\item حمایت از حزب‌الله لبنان (از ۱۹۸۲)
\item نفوذ در عراق پس از ۲۰۰۳
\item حمایت از دولت اسد در سوریه (از ۲۰۱۱)
\item حمایت از حوثی‌های یمن
\item شکل‌دهی \concept{محور مقاومت} (\lat{Axis of Resistance})
\end{itemize}
\end{tcolorbox}

نکته‌ی مهم آن است که ایران — بر خلاف عراق، لیبی، و افغانستان — به دلایل متعدد هدف مداخله‌ی نظامی مستقیم قرار نگرفته است: جمعیت بزرگ (۸۵ میلیون)، جغرافیای دشوار، ساختار نظامی نسبتاً قوی (شامل توانمندی موشکی)، حمایت روسیه و چین، و تجربه‌ی شکست مداخلات دیگر در منطقه. اما این به معنای عدم مداخله نیست: \concept{مداخله‌ی نرم} (تحریم، جنگ سایبری، عملیات رسانه‌ای) شکل غالب مداخله‌ی خارجی در ایران معاصر بوده است. این موضوع در پیوست~پ مفصلاً بحث خواهد شد.


%% ================================================================
\section{جمع‌بندی فصل}
\label{sec:ch10-conclusion}
%% ================================================================

\needspace{6\baselineskip}
فصل حاضر نشان داد که مداخلات معاصر — علی‌رغم تنوع در قالب، توجیه، و بازیگران — الگوهای مشترک بنیادینی دارند. از جنگ خلیج فارس ۱۹۹۱ تا تهاجم روسیه به اوکراین در ۲۰۲۲، هر مداخله‌ای حکایتی از تنش میان \concept{آرزواندیشی} و \concept{واقع‌گرایی}، میان \concept{شعار} و \concept{عمل}، و میان \concept{منافع} و \concept{ارزش‌ها} بوده است.

\begin{tcolorbox}[wavebox, title={یافته‌های کلیدی فصل ۱۰}]
\begin{enumerate}
\item \textbf{مشروعیت حقوقی لازم است اما کافی نیست:} حتی مداخلات دارای مجوز شورای امنیت (سومالی، لیبی) ممکن است شکست بخورند. مشروعیت حقوقی بدون برنامه‌ریزی عملی، ظرفیت اجرایی، و تعهد بلندمدت بی‌معناست.

\item \textbf{چرخش مأموریت قاعده است، نه استثنا:} در اکثریت مداخلات بررسی‌شده، اهداف اولیه به‌تدریج گسترش یافته و از هدف اصلی فاصله گرفته‌اند.

\item \textbf{عامل داخلی تعیین‌کننده‌تر از عامل خارجی است:} مقاومت اوکراین (۲۰۲۲) و ساختارهای مدنی کوزوو نشان دادند که بدون عامل داخلی قوی، هیچ مداخله‌ی خارجی نمی‌تواند ثبات پایدار ایجاد کند. برعکس، در عراق و لیبی، ضعف عامل داخلی به فروپاشی انجامید.

\item \textbf{اثر بومرنگ واقعی است:} مداخلات خارجی اغلب مشکلات جدیدی ایجاد می‌کنند که بدتر از مشکل اصلی‌اند: عراق → داعش، لیبی → بازار برده، افغانستان → بازگشت طالبان.

\item \textbf{انتخابی بودن مداخله مشروعیت کل نظام را زیر سؤال می‌برد:} وقتی R2P در لیبی اعمال شود اما در سوریه، یمن، و فلسطین نه، پرسش بنیادین این است: حقوق بشر حق است یا امتیاز؟

\item \textbf{فناوری ماهیت مداخله را تغییر می‌دهد:} جنگ سایبری، پهپادها، و شبکه‌های اجتماعی ابزارهای جدیدی برای مداخله‌ی بدون حضور نظامی مستقیم فراهم کرده‌اند — اما پرسش‌های اخلاقی و حقوقی جدیدی نیز ایجاد نموده‌اند.

\item \textbf{بازگشت رقابت قدرت‌های بزرگ:} اوکراین ۲۰۲۲ نشان داد که دوره‌ی مداخلات «آسان» غرب در کشورهای ضعیف (عراق، لیبی) پایان یافته و دوره‌ی جدید رقابت قدرت‌های بزرگ (آمریکا-چین-روسیه) آغاز شده است.

\item \textbf{نجات‌بخشی خارجی اسطوره‌ای خطرناک است:} تاریخ پنج دهه‌ی اخیر نشان می‌دهد که «نجات‌بخش خارجی» — چه آمریکا، چه روسیه، چه ناتو، چه عربستان — معمولاً منافع خود را دنبال می‌کند و ملت‌های «نجات‌یافته» اغلب در وضعیت بدتری قرار می‌گیرند. این یافته‌ی محوری، پلی است به فصل‌های ۱۱ و ۱۲ درباره‌ی ابعاد روان‌شناختی و جامعه‌شناختی نجات‌طلبی.
\end{enumerate}
\end{tcolorbox}

\begin{tcolorbox}[quotebox, title={سخن پایانی فصل}]
«ملت‌هایی که آزادی‌شان را مدیون بیگانگان‌اند، آزادی‌شان را حفظ نخواهند کرد. زیرا آنچه بدون هزینه و تلاش خودی به دست آید، بدون هزینه و تلاش از دست خواهد رفت.»\\
— مضمونی آزاد از \thinker{جان استوارت میل} (\lat{John Stuart Mill})، \textit{درباره‌ی آزادی}، ۱۸۵۹\footnote{\lr{Mill, John Stuart. "A Few Words on Non-Intervention," \textit{Fraser's Magazine}, December 1859. Reprinted in \textit{Dissertations and Discussions}, Vol. III, London: Longmans, Green, 1867, pp. 153–178.}}
\end{tcolorbox}

\cleardoublepage